\documentclass{article}

% Preamble
\usepackage[margin=1in]{geometry}     % For margins
\usepackage{amsmath, amssymb}         % For mathematics
\usepackage{graphicx}                 % For images (not used in this specific content but required)
\usepackage{hyperref}                 % For clickable links
\usepackage{xcolor}                   % For colors (used by hyperref, good practice)
\usepackage{booktabs}                 % For high quality professional tables
\usepackage{tabularx}                 % For tables with adjustable column width
\usepackage{listings}                 % For code blocks (not used in this specific content but required)
\usepackage{enumitem}                 % For customizable list environments
\usepackage{ragged2e}                 % For text alignment (e.g., \justifying)
\usepackage{titlesec}                 % For section formatting if needed (used for subsections without numbering)

% Basic clean typography
\setlength{\parskip}{0.5\baselineskip} % Space between paragraphs
\setlength{\parindent}{0pt}             % No paragraph indentation
\raggedbottom                           % Prevents extra vertical space filling on pages

% Hyperref setup
\hypersetup{
    colorlinks=true,
    linkcolor=blue,
    urlcolor=blue,
    citecolor=blue
}

% Custom command for unnumbered subsections in the final interdisciplinary section
\titleformat{\subsection}[runin]{\normalfont\bfseries}{\thesubsection}{1em}{}[]
\titlespacing*{\subsection}{0pt}{1.5ex plus 1ex minus .2ex}{1em}

% Title information
\title{Question Bank: Entrepreneurship, Innovation, and Business Analytics}
\author{Expert LaTeX Typesetter}
\date{Based on Various Exam Sessions (2012-2025)}

\begin{document}

\maketitle
\tableofcontents
\newpage

\section{Entrepreneurship Development – Concept, Types, Theories and Process, Developing Entrepreneurial Competencies}

\textbf{Q.87 (DEC 2023) | Direct MCQ}
"Innovation is the specific instrument of entrepreneurship - the act that endows resources with a new capacity to create wealth" who said this?
\begin{enumerate}[label=(\arabic*)]
    \item Peter Drucker
    \item J. A. Schumpeter
    \item Frank H. Knight
    \item Frank W. Young
\end{enumerate}

\textbf{Q.125 (DEC 2023) | Match the Following}
Match List - I with List - II.
\begin{tabularx}{\linewidth}{@{} p{0.5\linewidth} X @{}}
    \toprule
    \textbf{List - I Name of Theories of Entrepreneurship} & \textbf{List - II Proponents} \\
    \midrule
    (A) Theory of Opportunities and Innovation. & (I) Frank W. Young \\
    (B) Theory of Change. & (II) Peter F. Drucker \\
    (C) Theory of Religious Beliefs. & (III) John H. Kunkel \\
    (D) Theory of Entrepreneurship Supply. & (IV) Max Weber \\
    \bottomrule
\end{tabularx}
\begin{enumerate}[label=(\arabic*)]
    \item (A)-(II), (B)-(I), (C)-(IV), (D)-(III)
    \item (A)-(I), (B)-(II), (C)-(III), (D)-(IV)
    \item (A)-(IV), (B)-(III), (C)-(II), (D)-(I)
    \item (A)-(II), (B)-(III), (C)-(I), (D)-(IV)
\end{enumerate}

\textbf{Q.47 (JUN 2023) | Direct MCQ}
Entrepreneurs in general are _______ obsessed and _______ balanced?
\begin{enumerate}[label=(\arabic*)]
    \item Opportunity, Leadership
    \item Innovation, Risk
    \item Creativity, Innovations
    \item Profit, Creativity
\end{enumerate}

\textbf{Q.39 (MAR 2023) | Direct MCQ}
'Entrepreneur is an innovator playing the role of a dynamic businessman adding material growth to economic development', is articulated by:
\begin{enumerate}[label=(\arabic*)]
    \item Porter (1990)
    \item Lichtenstein and Lyons (2010)
    \item Joseph Schumpeter (1883-1950)
    \item Shane (2008)
\end{enumerate}

\textbf{Q.76 (MAR 2023) | Match the Following}
Match List I with List II
\begin{tabularx}{\linewidth}{@{} p{0.45\linewidth} X @{}}
    \toprule
    \textbf{List I (Entrepreneurship Thinkers)} & \textbf{List II (Entrepreneurial Manifestations)} \\
    \midrule
    A. Baumol (2002) & I. Innovating firms are more profitable than the non-innovators \\
    B. Geroski et al. (1993) and Leiponen (2000) & II. Entrepreneur is a marriage broker between what is desirable from an economic point of view and what is profitable from technological point of view \\
    C. Timmons and Spinelli (2007) & III. Entrepreneurship is a way of thinking, reasoning and acting that is opportunity obsessed, holistic in approach and leadership balanced \\
    D. Burton Klein (1917) & IV. In free market economies, innovation is a necessity if a firm wants to survive \\
    \bottomrule
\end{tabularx}
\begin{enumerate}[label=(\arabic*)]
    \item A-IV, B-I, C-III, D-II
    \item A-II, B-IV, C-III, D-I
    \item A-I, B-IV, C-II, D-III
    \item A-III, B-I, C-IV, D-II
\end{enumerate}

\textbf{Q.91 (MAR 2023) | Assertion-Reasoning}
Given below are two statements:
Assertion (A): Entrepreneurship has emerged as a focal point of policy in the majority of countries.
Reason (R): Entrepreneurship contributes in job generation, innovation and promotes social cohesion.
\begin{enumerate}[label=(\arabic*)]
    \item Both (A) and (R) are correct and (R) is the correct explanation of (A).
    \item Both (A) and (R) are correct but (R) is NOT the correct explanation of (A).
    \item (A) is correct but (R) is incorrect.
    \item (A) is incorrect but (R) is correct.
\end{enumerate}

\textbf{Q.89 (2 SEP 2024) | Direct MCQ}
What can lead to the withdrawal of status respect according to Hagen's theory of Entrepreneurship?
\begin{enumerate}[label=(\arabic*)]
    \item Increased recognition from the group
    \item Perception that efforts and purposes are highly valued
    \item Belief in the stability of status symbols
    \item Perception that efforts and purposes are not valued by others
\end{enumerate}

\textbf{Q.140 (2 SEP 2024) | Match the Following}
Match the List-I with List-II
\begin{tabularx}{\linewidth}{@{} p{0.5\linewidth} X @{}}
    \toprule
    \textbf{LIST I (Name of Theories)} & \textbf{LIST II (Propounder)} \\
    \midrule
    A. Theory of Opportunities and Innovation & I. Frank H Knight \\
    B. Theory of profit & II. Peter F Drucker \\
    C. Theory of Entrepreneurial Supply & III. Sigmund Freud \\
    D. Psychoanalytic Theory & IV. Thomas Cochran \\
    \bottomrule
\end{tabularx}
\begin{enumerate}[label=(\arabic*)]
    \item A-II, B-III, C-IV, D-I
    \item A-II, B-I, C-IV, D-III
    \item A-I, B-III, C-II, D-IV
    \item A-III, B-II, C-I, D-IV
\end{enumerate}

\textbf{Q.89 (5 SEP 2024) | Direct MCQ}
Which of the following is not part of the common myths about entrepreneurs?
\begin{enumerate}[label=(\arabic*)]
    \item Entrepreneurs are gamblers
    \item Entrepreneurs are born, not made
    \item Entrepreneurs are motivated primarily by money
    \item Entrepreneurs are creative
\end{enumerate}

\textbf{Q.88 (JUN 2025) | Direct MCQ}
Which of the following is the example of Social Entrepreneurship ?
\begin{enumerate}[label=(\arabic*)]
    \item PFL
    \item Sulabh International
    \item Hyundai
    \item HUL
\end{enumerate}

\textbf{Q.90 (JUN 2025) | Direct MCQ}
Which of the following is not the type of start up firms ?
\begin{enumerate}[label=(\arabic*)]
    \item Salary-substitute firms
    \item Life style firms
    \item Entrepreneurial firms
    \item Ordinance firms
\end{enumerate}

\textbf{Q.124 (JUN 2025) | Multi-Statement}
Which of the following types of entrepreneurs are based on personality traits ?
\begin{enumerate}[label=\Alph*.]
    \item Franchises
    \item The real manager
    \item The achiever
    \item Industrial entrepreneur
    \item Owner-cum-manager
\end{enumerate}
\begin{enumerate}[label=(\arabic*)]
    \item A, B, C Only
    \item B and C Only
    \item B, C, D Only
    \item C and E Only
\end{enumerate}

\textbf{Q.54 (DEC 2025) | Chronological Sequence}
Identify the correct sequence of entrepreneurial process :
\begin{enumerate}[label=\Alph*.]
    \item Establish vision
    \item Identify an opportunity
    \item Gather resources
    \item Persuade others
    \item Create New venture/product/market
\end{enumerate}
\begin{enumerate}[label=(\arabic*)]
    \item B, A, D, C, E
    \item B, D, A, C, E
    \item A, B, D, C, E
    \item B, E, D, C, A
\end{enumerate}

\textbf{Q.90 (DEC 2025) | Direct MCQ}
An entrepreneur whose aim is to maximise the economic returns at a level consistent with the survival of the firm, with or without the element of growth is called :
\begin{enumerate}[label=(\arabic*)]
    \item Fabian entrepreneur
    \item First generation entrepreneur
    \item Modern entrepreneur
    \item Classical entrepreneur
\end{enumerate}

\textbf{Q.98 (DEC 2025) | Assertion-Reasoning}
Given below are two statements : one is labelled as Assertion (A) and the other is labelled as Reason (R).
Assertion (A) : Most fast growth entrepreneurial ventures organise as corporations or limited liability companies rather than as sole proprietorship or partnership.
Reason (R) : Business losses of corporations can be deducted against the shareholder's other sources of income.
\begin{enumerate}[label=(\arabic*)]
    \item Both (A) and (R) are correct and (R) is the correct explanation of (A)
    \item Both (A) and (R) are correct but (R) is not the correct explanation of (A)
    \item (A) is correct but (R) is not correct
    \item (A) is not correct but (R) is correct
\end{enumerate}

\textbf{Q.76 (DEC 2019) | Direct MCQ}
Which national institute has set up/established a 'Centre for Research in Entrepreneurship Education and Development (CREED)' with the objective to broaden the frontier of entrepreneurship research?
\begin{enumerate}[label=(\arabic*)]
    \item Indian Institute of Management, Ahmedabad
    \item Entrepreneurship Development Institute of India
    \item Delhi University
    \item Institute of Rural Management, Anand
\end{enumerate}

\textbf{Q.108 (4 DEC 2019) | Multi-Statement}
Entrepreneurs are motivated to start business enterprise because of various factors :
\begin{enumerate}[label=\alph*.]
    \item Ambitious factors
    \item Compelling factors
    \item Facilitating factors
    \item Family factors
\end{enumerate}
Which of the following options is correct?
\begin{enumerate}[label=(\arabic*)]
    \item (a), (b) and (c) only
    \item (a) and (c) only
    \item (a) and (d) only
    \item (a), (b), (c) and (d)
\end{enumerate}

\textbf{Q.141 (24 JUN 2019) | Direct MCQ}
Which national institute has set up/established a 'Centre for Research in Entrepreneurship Education and Development (CREED)' with the objective to broaden the frontier of entrepreneurship research?
\begin{enumerate}[label=(\arabic*)]
    \item Indian Institute of Management, Ahmedabad
    \item Entrepreneurship Development Institute of India
    \item Delhi University
    \item Institute of Rural Management, Anand
\end{enumerate}

\textbf{Q.115 (SEP 2020) | Multi-Statement}
The statement "Entrepreneurship is a mother of foresight and willingness to assume risks in undertakings particularly a new one". Who has said above the underlined portion of the statement?
\begin{enumerate}[label=\Alph*.]
    \item Richard Cantillon.
    \item Leon Walras.
    \item William Diamond.
    \item Peter F. Drucker.
    \item J.A. Schumpeter.
\end{enumerate}
Choose the correct answer from the options given below:
\begin{enumerate}[label=(\arabic*)]
    \item A, C only
    \item B, C only
    \item D, E only
    \item C, E only
\end{enumerate}

\textbf{Q.39 (NOV 2021) | Direct MCQ}
Entrepreneurs who are characterised by a refusal to adopt opportunities and make changes in production systems even if it affects the returns of the organisation are called
\begin{enumerate}[label=(\arabic*)]
    \item Fabian Entrepreneurs
    \item Life timers entrepreneurs
    \item Imitative entrepreneurs
    \item Drone entrepreneurs
\end{enumerate}

\textbf{Q.51 (NOV 2021) | Multi-Statement MCQ}
Out of the following, which are the most common error committed by Entrepreneurs?
\begin{enumerate}[label=\Alph*.]
    \item To anticipate roadblocks
    \item No commitment or dedication
    \item Proposing market niche (Segment)
    \item Lack of demonstrated experience
    \item No realistic goals
\end{enumerate}
\begin{enumerate}[label=(\arabic*)]
    \item A, B and C only
    \item B, D and E only
    \item C, D and E only
    \item B, C, D and E only
\end{enumerate}

\textbf{Q.53 (NOV 2021) | Multi-Statement MCQ}
Notable myths about Entrepreneurs and Entrepreneurship are
\begin{enumerate}[label=\Alph*.]
    \item Entrepreneurs are academics and social misfits
    \item Entrepreneurs must fit into an ideal profile
    \item Entrepreneurs are doers and thinkers
    \item All Entrepreneurs need is money
    \item Entrepreneurs are not born but made
\end{enumerate}
\begin{enumerate}[label=(\arabic*)]
    \item A, B, C and D only
    \item B, C, D and E only
    \item A, B and D only
    \item C, D and E only
\end{enumerate}

\textbf{Q.75 (NOV 2021) | Match the Following}
Match List I with List II
\begin{tabularx}{\linewidth}{@{} p{0.45\linewidth} X @{}}
    \toprule
    \textbf{List I (Entrepreneur type)} & \textbf{List II (Features based on Entrepreneurial Competency)} \\
    \midrule
    A. Passenger & I. Ability to do and no will to do \\
    B. Star/Ideal & II. No ability to do and no will to do \\
    C. Dead wood & III. Ability to do and will to do \\
    D. May run & IV. No ability to do and will to do \\
    \bottomrule
\end{tabularx}
\begin{enumerate}[label=(\arabic*)]
    \item A-I, B-II, C-III, D-IV
    \item A-IV, B-III, C-II, D-I
    \item A-IV, B-III, C-I, D-II
    \item A-III, B-IV, C-I, D-II
\end{enumerate}

\textbf{Q.38 (2022) | Direct MCQ}
Currently, the country with the highest level of entrepreneurial activity, represented by Global Entrepreneurship Index is:
\begin{enumerate}[label=(\arabic*)]
    \item Korea
    \item Norway
    \item US
    \item Australia
\end{enumerate}

\textbf{Q.118 (DEC 2015) | Assertion-Reasoning}
Assertion (A) : The function of entrepreneurs is to reform or revolutionise the pattern of production by exploiting an invention or, more generally, an untried technological possibility for producing a new commodity or producing an old one in a new way.
Reasoning (R) : Lack of innovations will usually reduce the life span of a product and/or service.
Code :
\begin{enumerate}[label=(\arabic*)]
    \item Assertion (A) and Reasoning (R), both, are correct, and (R) is correct explanation of (A).
    \item Assertion (A) and Reasoning (R), both, are correct, but (R) is not right explanation of (A).
    \item Assertion (A) is correct, whereas Reasoning (R) is incorrect.
    \item Reasoning (R) is correct, whereas Assertion (A) is incorrect.
\end{enumerate}

\textbf{Q.120 (DEC 2015) | Multi-Statement}
Which of the following are correct for an entrepreneur ?
\begin{enumerate}[label=(\alph*)]
    \item is an independent person
    \item makes his/her own decisions
    \item has the opportunity of realising dreams
    \item has good financial strength
    \item has undergone a systematic business training
    \item has immense job satisfaction
\end{enumerate}
Code :
\begin{enumerate}[label=(\arabic*)]
    \item (a), (b), (c) and (d)
    \item (a), (b), (c) and (e)
    \item (a), (b), (c) and (f)
    \item (a), (b), (d) and (f)
\end{enumerate}

\textbf{Q.124 (DEC 2015) | Match the Following}
Match the items of List – I with List – II with regards to phases of entrepreneurial process :
\begin{tabularx}{\linewidth}{@{} l X @{}}
    \toprule
    \textbf{List - I} & \textbf{List - II} \\
    \midrule
    (a) Identify and evaluate the opportunity & (i) Executive summary \\
    (b) Develop business plan & (ii) Develop access to needed resources \\
    (c) Determine the required resources & (iii) Implement control system \\
    (d) Manage the enterprise & (iv) Assess the competitive environment \\
    \bottomrule
\end{tabularx}
Codes:
\begin{enumerate}[label=(\arabic*)]
    \item (a)-(iv), (b)-(ii), (c)-(iii), (d)-(i)
    \item (a)-(iv), (b)-(i), (c)-(ii), (d)-(iii)
    \item (a)-(i), (b)-(iii), (c)-(ii), (d)-(iv)
    \item (a)-(iv), (b)-(iii), (c)-(i), (d)-(ii)
\end{enumerate}

\textbf{Q.40 (AUG 2016 P-III) | Multi-Statement}
Consider the following statements :
Statement - I: The entrepreneur has to be a mere manager. He (or she) may not be an innovator and promoter as well.
Statement – II : An entrepreneur always searches for change, responds to it and exploits it as an opportunity.
Statement – III : The government should not try to push the entrepreneurship beyond the boundaries of a business unit.
Code:
\begin{enumerate}[label=(\arabic*)]
    \item Statement I is only true, others are false.
    \item Statement II is only true, others are false.
    \item Statement II \& III are true, statement I is false.
    \item All the statements are true.
\end{enumerate}

\textbf{Q.45 (AUG 2016 P-II) | Match the Following}
Match the items of List - I with those of List – II.
\begin{tabularx}{\linewidth}{@{} p{0.4\linewidth} X @{}}
    \toprule
    \textbf{List - I (Entrepreneurship Theory)} & \textbf{List - II (Attributes)} \\
    \midrule
    (a) Social change & (i) The introduction of a new product which consumers are not yet familiar. \\
    (b) Status withdrawal theory & (ii) Individuals with high need achievement will not be motivated by monetary incentives but that monetary rewards will constitute a symbol of achievement for them. \\
    (c) Achievement theory & (iii) Driving entrepreneurial energies are generated by the adoption of exogenously supplied religious beliefs. \\
    (d) Innovation theory & (iv) A class which lost its previous prestige or a minority group tends to show aggressive entrepreneurial drive. \\
    \bottomrule
\end{tabularx}
Codes:
\begin{tabularx}{\linewidth}{@{} l *{4}{c} @{}}
    \toprule
    & (a)   & (b)   & (c)  & (d) \\
    \midrule
    (1)   & (iii) & (i)  & (ii) & (iv) \\
    (2)   & (iii) & (iv) & (ii) & (i) \\
    (3)   & (i)   & (ii) & (iii) & (iv) \\
    (4)   & (iv)  & (ii) & (iii) & (i) \\
    \bottomrule
\end{tabularx}

\textbf{Q.41 (JAN 2017 P-II) | Match the Following}
Match the items of List – I with the items of List – II and select the code of correcting matching:
\begin{tabularx}{\linewidth}{@{} p{0.3\linewidth} X @{}}
    \toprule
    \textbf{List - I} & \textbf{List - II} \\
    \midrule
    a. Richard Cantillon         & i. Entrepreneur must accept the challenge and should be willing hard to achieve something \\
    b. Joseph A. Schumpeter      & ii. Entrepreneurs are all those persons who are engaged in economic activity. \\
    c. J.K. Galbraith            & iii. Entrepreneur is a person who introduces innovation and change. \\
    \bottomrule
\end{tabularx}
\begin{tabularx}{\linewidth}{@{} l *{3}{c} @{}}
    \toprule
    & a    & b    & c \\
    \midrule
    (1)  & iii  & i    & ii \\
    (2)  & i    & iii  & ii \\
    (3)  & ii   & iii  & i \\
    (4)  & ii   & i    & iii \\
    \bottomrule
\end{tabularx}

\textbf{Q.42 (JAN 2017 P-II) | Direct MCQ}
Total Quality Management, job redesigning, new techniques of doing things and management by consensus are the examples of which one of the following?
\begin{enumerate}[label=(\arabic*)]
    \item Opportunistic entrepreneurship
    \item Administrative entrepreneurship
    \item Incubative entrepreneurship
    \item Mass entrepreneurship
\end{enumerate}

\textbf{Q.44 (JAN 2017 P-II) | Direct MCQ}
Which one of the following is not a personal characteristic of an entrepreneur?
\begin{enumerate}[label=(\arabic*)]
    \item Innovative
    \item Hardworking
    \item Practical
    \item Submissive
\end{enumerate}

\textbf{Q.33 (JAN 2017 P-III) | Direct MCQ}
Which one of the following is not a common motivation of an entrepreneur?
\begin{enumerate}[label=(\arabic*)]
    \item Desire to maintain the status quo
    \item Search for more rewarding work
    \item Search for personal and professional growth
    \item Desire for independence
\end{enumerate}

\textbf{Q.35 (JAN 2017 P-III) | Match the Following}
Match the items of List – I with the items of List – II and identify the correct code of combination:
\begin{tabularx}{\linewidth}{@{} p{0.45\linewidth} X @{}}
    \toprule
    \textbf{List - I} & \textbf{List - II} \\
    \midrule
    a. Entrepreneurship and Dynamic Capitalism   & i. A.D. Chandler Jr. \\
    b. The Affluent Society                     & ii. J.K. Galbraith \\
    c. Theory of Economic Development           & iii. Joseph Schumpeter \\
    d. The Dynamics of Industrial Capitalism    & iv. B.A. Kirchhoff \\
    \bottomrule
\end{tabularx}
\begin{tabularx}{\linewidth}{@{} l *{4}{c} @{}}
    \toprule
    & a    & b    & c    & d \\
    \midrule
    (1)  & i    & iii  & ii   & iv \\
    (2)  & i    & ii   & iii  & iv \\
    (3)  & iv   & iii  & ii   & i \\
    (4)  & iv   & ii   & iii  & i \\
    \bottomrule
\end{tabularx}

\textbf{Q.36 (JAN 2017 P-III) | Statement-I \& II}
Statement – I: Entrepreneurs recognize opportunities where others see chaos or confusion.
Statement-II: Entrepreneurs bring aggressive capital for the existing business in traditional market place.
\begin{enumerate}[label=(\arabic*)]
    \item Statement – I is correct but II is not correct.
    \item Statement – II is correct but I is not correct.
    \item Both the Statements are correct.
    \item Both the Statements are incorrect.
\end{enumerate}

\textbf{Q.41 (NOV 2017 P-II) | Direct MCQ}
Elements of entrepreneurial competencies, according to the theorist of the subject, pertain to which one of the following?
\begin{enumerate}[label=(\arabic*)]
    \item Body of knowledge, set of skills, cluster of appropriate motives and traits
    \item Ability to prepare viable project, executional skills and materialisation of group goals
    \item Ability to work in a group, enforce group dynamics and motivate employees
    \item Assigning priorities to hierarchy of Maslow's monetary needs and non-monetary needs
\end{enumerate}

\textbf{Q.48 (DEC 2018) | Match the Following}
Match the items of List I with the items of List II and choose the correct answer from the code given below.
\begin{tabularx}{\linewidth}{@{} p{0.4\linewidth} X @{}}
    \toprule
    \textbf{List I} & \textbf{List II} \\
    \midrule
    (a) Fabian Entrepreneurs & (i) Industry leaders \\
    (b) Drone Entrepreneurs & (ii) Found in the places where there is a lack of resources \\
    (c) Imitating Entrepreneurs & (iii) Reluctant to change \\
    (d) Innovative Entrepreneurs & (iv) Sceptical about the changes to be made in the organization \\
    \bottomrule
\end{tabularx}
Codes:
\begin{enumerate}[label=(\arabic*)]
    \item (a)-(iv), (b)-(ii), (c)-(iii), (d)-(i)
    \item (a)-(iv), (b)-(iii), (c)-(ii), (d)-(i)
    \item (a)-(iii), (b)-(ii), (c)-(iv), (d)-(i)
    \item (a)-(i), (b)-(ii), (c)-(iv), (d)-(iii)
\end{enumerate}

\textbf{Q.66 (DEC 2018) | Multi-Statement}
Who among the following are associated with the concept of "Entrepreneurship"?
\begin{enumerate}[label=\arabic*.]
    \item Richard Cantillon
    \item Lionel Robbins
    \item Adam Smith
    \item Jean-Baptiste Say
    \item Joseph Schumpeter
\end{enumerate}
Choose the correct answer from the code given below:
\begin{enumerate}[label=(\arabic*)]
    \item (i), (ii), (iv) and (v)
    \item (i), (iii), (iv) and (v)
    \item (ii), (iii), (iv) and (v)
    \item (i), (ii), (iii) and (iv)
\end{enumerate}

\textbf{Q.58 (JULY 2018) | Direct MCQ}
Which one among the following is not a factor external to an entrepreneur ?
\begin{enumerate}[label=(\arabic*)]
    \item Machinery on hire purchase
    \item Accommodation in industrial estates
    \item Financial assistance from non-government sources
    \item Business experience in the same or related line
\end{enumerate}

\textbf{Q.60 (JULY 2018) | Direct MCQ}
Individuals who are the founders of the business and conceptualize a business plan by putting efforts to make the plan a success are known as :
\begin{enumerate}[label=(\arabic*)]
    \item Quasi entrepreneurs
    \item Bold entrepreneurs
    \item Pure entrepreneurs
    \item Owner-manager
\end{enumerate}

\textbf{Q.61 (JULY 2018) | Match the Following}
Match the items of List (II) with the items of List (I) :
\begin{tabularx}{\linewidth}{@{} p{0.55\linewidth} X @{}}
    \toprule
    \textbf{List (I)} & \textbf{List (II)} \\
    \midrule
    (a) Entrepreneurship is essentially a creative activity. & (i) Fransis A. Walker \\
    (b) An entrepreneur is an organizer and coordinator of various factors of production. & (ii) Peter F. Drucker \\
    (c) An entrepreneur is one who always searches for change, responds to it, and exploits it as an opportunity. & (iii) Joseph Schumpeter \\
    \bottomrule
\end{tabularx}
Code :
\begin{enumerate}[label=(\arabic*)]
    \item (a)-(iii), (b)-(i), (c)-(ii)
    \item (a)-(iii), (b)-(ii), (c)-(i)
    \item (a)-(ii), (b)-(i), (c)-(iii)
    \item (a)-(ii), (b)-(iii), (c)-(i)
\end{enumerate}

\textbf{Q.63 (JULY 2018) | Direct MCQ}
Who among the following bridges the gap between inventors and managers ?
\begin{enumerate}[label=(\arabic*)]
    \item Investor
    \item Entrepreneur
    \item Industrialist
    \item Intrapreneur
\end{enumerate}

\textbf{Q.48 (Archived Set 1) | Direct MCQ}
“Innovation and entrepreneurship” is a book written by
\begin{enumerate}[label=(\arabic*)]
    \item Philip Kotler
    \item C. K. Prahlad
    \item Pradeep Khandwala
    \item Peter F. Drucker
\end{enumerate}

\textbf{Q.PB-44 (Archived MSME Policies Set) | Direct MCQ}
EDII is situated at
\begin{enumerate}[label=(\arabic*)]
    \item New Delhi
    \item Ahmedabad
    \item Patna
    \item Kolkata
\end{enumerate}

\textbf{Q.42 (DEC 2012 P-III) | Direct MCQ}
People who take the risks necessary to organize, manage and receive the financial profits and non-monetary rewards are called
\begin{enumerate}[label=(\arabic*)]
    \item Suppliers
    \item Employees
    \item Competitors
    \item Entrepreneurs
\end{enumerate}

\textbf{Q.41 (DEC 2012 P-II) | Direct MCQ}
Entrepreneurship as a theory of business was propounded by
\begin{enumerate}[label=(\arabic*)]
    \item Douglas McGregor
    \item Thomas A. Coleman
    \item Joseph A. Schumpeter
    \item Adam Smith
\end{enumerate}

\textbf{Q.43 (DEC 2012 P-II) | Direct MCQ}
The book ‘Innovation and Entrepreneurship’ is written by
\begin{enumerate}[label=(\arabic*)]
    \item Harold Koontz
    \item Blake and Mouton
    \item Peter F. Drucker
    \item None of the above
\end{enumerate}

\textbf{Q.41 (JUNE 2012) | Direct MCQ}
Which one is not the characteristic feature of Entrepreneurship ?
\begin{enumerate}[label=(\arabic*)]
    \item Vision
    \item Risk Bearing
    \item Initiative and Drive
    \item Disloyalty
\end{enumerate}

\textbf{Q.86 (JUNE 2012) | Direct MCQ}
Activities taken up on part time or casual basis to raise income is :
\begin{enumerate}[label=(\arabic*)]
    \item Self employment
    \item Income generation
    \item Entrepreneurship
    \item None of the above
\end{enumerate}

\textbf{Q.41 (DEC 2013) | Direct MCQ}
Who said – “An entrepreneur always searches for change, responds to it and exploit it as an opportunity.” ?
\begin{enumerate}[label=(\arabic*)]
    \item James Burna
    \item McClelland
    \item Peter F. Drucker
    \item Robert C. Ronstadt
\end{enumerate}

\textbf{Q.42 (DEC 2013) | Direct MCQ}
‘Kakinada Experiment’ on achievement motivation was conducted in
\begin{enumerate}[label=(\arabic*)]
    \item Uttar Pradesh
    \item Jammu \& Kashmir
    \item Andhra Pradesh
    \item Gujarat
\end{enumerate}

\textbf{Q.43 (DEC 2013) | Chronological Sequence}
The entrepreneurial process consists of the following steps :
\begin{enumerate}[label=\roman*.]
    \item Management of the resulting enterprise.
    \item Determination of the required resources.
    \item Identification and evaluation of the opportunity.
    \item Development of the business plan.
\end{enumerate}
Indicate the correct sequence.
\begin{enumerate}[label=(\arabic*)]
    \item iii iv i ii
    \item iv i iii ii
    \item iii iv ii i
    \item iv i ii iii
\end{enumerate}

\textbf{Q.108 (DEC 2013) | Direct MCQ}
An entrepreneur who is neither willing to introduce new changes nor to adopt new methods is known as :
\begin{enumerate}[label=(\arabic*)]
    \item Adoptive Entrepreneur
    \item Fabian Entrepreneur
    \item Innovative Entrepreneur
    \item Drone Entrepreneur
\end{enumerate}

\textbf{Q.116 (JUNE 2013) | Direct MCQ}
"An entrepreneur always searches for change, responds to it and exploit it as an opportunity.” Who said it ?
\begin{enumerate}[label=(\arabic*)]
    \item F.W. Taylor
    \item Peter F. Drucker
    \item J.R. Tulsian
    \item Max Weber
\end{enumerate}

\textbf{Q.118 (JUNE 2013) | Direct MCQ}
Which is not the object of Entrepreneurship Development Programmes ?
\begin{enumerate}[label=(\arabic*)]
    \item To create successful entrepreneur
    \item To remove doubts of entrepreneurs and to give solutions to the problems
    \item To create awareness about Government schemes and programmes
    \item None of the above
\end{enumerate}

\textbf{Q.P2-Q41 (SEP 2013) | Direct MCQ}
An Entrepreneur who is dominated more by customs, religions, traditions and past practices and he is not ready to take any risk is called as
\begin{enumerate}[label=(\arabic*)]
    \item Drone Entrepreneur
    \item Adoptive Entrepreneur
    \item Fabian Entrepreneur
    \item Innovative Entrepreneur
\end{enumerate}

\textbf{Q.P2-Q42 (SEP 2013) | Direct MCQ}
“Poverty is an artificial creation. It does not belong to human civilization and we can change that and can make people come out of poverty through redesigning our institutions and policies.” This preamble refers to
\begin{enumerate}[label=(\arabic*)]
    \item Women Entrepreneurship
    \item Corporate Entrepreneurship
    \item Social Entrepreneurship
    \item None of the above
\end{enumerate}

\textbf{Q.P3-Q56 (SEP 2013) | Chronological Sequence}
Arrange the Entrepreneurial motivation factors on the basis of proper sequences.
\begin{enumerate}[label=(\arabic*)]
    \item Need for achievement, Locus of control, vision, desire for independence, passion and drive.
    \item Locus of control, vision, desire for independence, passion, drive and need for achievement.
    \item Vision, desire for independence, passion, drive, locus of control and need for achievement.
    \item Desire for independence, need for achievement, locus of control, vision, passion and drive.
\end{enumerate}

\textbf{Q.P2-Q43 (SEP 2013) | Match the Following}
Match the following Management subjects with their techniques (relating to Entrepreneurship and Innovation processes).
\emph{Note: List I, List II and options for this Match the Following question are missing in the provided Markdown.}

\textbf{Q.41 (DEC 2014) | Direct MCQ}
Who propounded entrepreneurship as a theory of business ?
\begin{enumerate}[label=(\arabic*)]
    \item Daniel Goleman
    \item Joseph A. Schumpeter
    \item Thomas A. Coleman
    \item Warren Buffet
\end{enumerate}

\textbf{Q.42 (DEC 2014) | Direct MCQ}
Who among the following is not the first generation entrepreneur ?
\begin{enumerate}[label=(\arabic*)]
    \item Dhirubhai Ambani
    \item Shiv Nadar
    \item K. Anji Reddy
    \item G.D. Birla
\end{enumerate}

\textbf{Q.43 (DEC 2014) | Direct MCQ}
Which of the following is not a possible route to market entry for a small business owner/entrepreneur ?
\begin{enumerate}[label=(\arabic*)]
    \item Franchise
    \item Corporate Venture
    \item Outright Purchase
    \item Buy-out
\end{enumerate}

\textbf{Q.109 (DEC 2014) | Direct MCQ}
Which of the following is not a misconception about entrepreneurship ?
\begin{enumerate}[label=(\arabic*)]
    \item ‘Successful entrepreneurship needs only a great idea.’
    \item ‘Entrepreneurship is an easy task.’
    \item ‘Entrepreneurship is found only in small business.’
    \item ‘Entrepreneurial ventures and small business are two different things.’
\end{enumerate}

\textbf{Q.111 (DEC 2014) | Direct MCQ}
In which type of entrepreneurship formal structural ties are loosened and champions are given the freedom to pursue opportunities both for the organization and through market ?
\begin{enumerate}[label=(\arabic*)]
    \item Administrative
    \item Acquisitive
    \item Opportunistic
    \item Incubative
\end{enumerate}

\textbf{P-II Q.41 (JUNE 2014) | Direct MCQ}
A study conducted by the Entrepreneurship Development Institute of India, Ahmedabad, revealed that the possession of competencies is necessary for superior performance of the entrepreneurs. The study was conducted under the guidance of
\begin{enumerate}[label=(\arabic*)]
    \item James T. McCrory
    \item David C. McClelland
    \item James J. Berna
    \item P.L. Tandon
\end{enumerate}

\textbf{P-III Q.60 (JUNE 2014) | Match the Following}
Match the following Types of Qualities (Physical, Ethical, Social, Mental) with Features (Sharp memory, Coordinator, Personality, Character).
\emph{Note: List I, List II and options for this Match the Following question are missing in the provided Markdown.}

\section{Intrapreneurship – Concept and Process}

\textbf{Q.89 (DEC 2023) | Direct MCQ}
Which type of Intrapreneur is known for generating innovative ideas and seeking ways to improve processes?
\begin{enumerate}[label=(\arabic*)]
    \item Employee Intrapreneur
    \item Doers Intrapreneur
    \item Creator Intrapreneur
    \item Implementers Intrapreneur
\end{enumerate}

\textbf{Q.87 (2 SEP 2024) | Direct MCQ}
Which type of Intrapreneur is known for generating innovative ideas and seeking ways to improve process?
\begin{enumerate}[label=(\arabic*)]
    \item Employee Intrapreneur
    \item Doers Intrapreneur
    \item Creator Intrapreneur
    \item Implementers Intrapreneur
\end{enumerate}

\textbf{Q.87 (JUN 2025) | Direct MCQ}
Which of the following is not essential for nurturing Intrapreneurs ?
\begin{enumerate}[label=(\arabic*)]
    \item Building culture and climate
    \item Building Competence
    \item Building relationship
    \item Learning Sign language
\end{enumerate}

\textbf{Q.7 (DEC 2019) | Direct MCQ}
Which among the following is a correct statement?
\begin{enumerate}[label=(\arabic*)]
    \item Entrepreneur is an employee and intrapreneur is free and the leader of the operation.
    \item Intrapreneur is an employee and entrepreneur is free and the leader of the operation.
    \item Entrepreneurship is the change initiative taken within a going concern by the people working in the organisation.
    \item Intrapreneurship is the change initiative taken by a group of independent people.
\end{enumerate}

\textbf{Q.143 (24 JUN 2019) | Direct MCQ}
Which among the following is a correct statement?
\begin{enumerate}[label=(\arabic*)]
    \item Entrepreneur is an employee and intrapreneur is free and the leader of the operation.
    \item Intrapreneur is an employee and entrepreneur is free and the leader of the operation.
    \item Entrepreneurship is the change initiative taken within a going concern by the people working in the organisation.
    \item Intrapreneurship is the change initiative taken by a group of independent people.
\end{enumerate}

\textbf{Q.28 (DEC 2018) | Two-Statement}
Choose the correct code for the following statements being correct or incorrect.
Statement I: An Intrapreneur is a person within a large corporation who takes direct responsibility for turning an idea into a profitable finished product through assertive risk-taking and innovation.
Statement II: Intrapreneurship refers to employee initiatives in organizations to undertake something new, without being asked to do so.
Codes:
\begin{enumerate}[label=(\arabic*)]
    \item Statement I is correct but II is incorrect.
    \item Both the statements I and II are correct.
    \item Both the statements I and II are incorrect.
    \item Statement II is correct, but I is incorrect.
\end{enumerate}

\textbf{Q.17 (DEC 2012 P-III) | Direct MCQ}
An entrepreneurial person employed by a corporation and encouraged to be innovative and creative is referred to as
\begin{enumerate}[label=(\arabic*)]
    \item Competitor
    \item Supplier
    \item Entrepreneur
    \item Intrapreneur
\end{enumerate}

\textbf{Q.110 (DEC 2014) | Direct MCQ}
Pichot suggested three roles for the intrapreneur in promoting innovation. Indicate the correct answer.
\begin{enumerate}[label=(\arabic*)]
    \item Inventor, innovator and sponsor
    \item Inventor, product champion and sponsor
    \item Team builder, information source and advocate
    \item Mould breaker, change agent and sponsor
\end{enumerate}

\section{Women Entrepreneurship and Rural Entrepreneurship}

\textbf{Q.88 (DEC 2023) | Direct MCQ}
In which institutions of women entrepreneurship was Jyoti Jeevan Naik associated?
\begin{enumerate}[label=(\arabic*)]
    \item Shri Mahila Griha Udyog Lijjat Papad
    \item Sahkari Mahila Gram Samittee
    \item Jivika
    \item Self-employed Women's Association
\end{enumerate}

\textbf{Q.26 (JUN 2023) | Multi-Statement}
The major problems faced by rural entrepreneurs are:
\begin{enumerate}[label=\alph*.]
    \item Rate of interest charged by commercial banks is too high
    \item Inability to market their products and services
    \item Lack of comprehensive training on required skills
    \item Inability to handle cultural issues
    \item Difficulty in opening bank accounts due to lack of documents
\end{enumerate}
\begin{enumerate}[label=(\arabic*)]
    \item a and b only
    \item a, d and e only
    \item b and c only
    \item c, d and e only
\end{enumerate}

\textbf{Q.27 (JUN 2023) | Assertion-Reasoning}
Assertion ‘A’: A stereotyped and male-centered vision discourages some women from participating in business activities that could have a consequence on people who interact with women at the community level.
Reason ‘R’: Attributions of the society and the different socialisation processes relating to men and women may create obstacles for women due to unequal distribution of assets and services.
\begin{enumerate}[label=(\arabic*)]
    \item Both A and R are true and R is the correct explanation of A
    \item Both A and R are true but R is NOT the correct explanation of A
    \item A is true but R is false
    \item A is false but R is true
\end{enumerate}

\textbf{Q.88 (2 SEP 2024) | Direct MCQ}
In which of the following Association/organisation Ela Bhatt is associated?
\begin{enumerate}[label=(\arabic*)]
    \item Mukta Services
    \item Lijjat Papad
    \item Self employed women's Association
    \item Avya Global connect
\end{enumerate}

\textbf{Q.140 (JUN 2025) | Match the Following}
Match List - I with List - II :
\begin{tabularx}{\linewidth}{@{} p{0.4\linewidth} X @{}}
    \toprule
    \textbf{List - I (Women Entrepreneur)} & \textbf{List - II (Enterprise)} \\
    \midrule
    A. Richa Kar & I. Zivame \\
    B. Kiran Mazumdar Shaw & II. Kamani Tubes \\
    C. Kalpana Saroj & III. Biocon Limited \\
    D. Falguni Nayar & IV. Nykaa \\
    \bottomrule
\end{tabularx}
\begin{enumerate}[label=(\arabic*)]
    \item A-I, B-III, C-II, D-IV
    \item A-I, B-II, C-III, D-IV
    \item A-I, B-III, C-IV, D-II
    \item A-I, B-IV, C-II, D-III
\end{enumerate}

\textbf{Q.69 (DEC 2019) | Direct MCQ}
NITI Aayog, GOI, has launched a scheme to develop ecosystem for women entrepreneurship in India. Identify the scheme out of the following:
\begin{enumerate}[label=(\arabic*)]
    \item Women Empowerment Agency
    \item Women Entrepreneurship Platform
    \item Naari Shakti Aayog
    \item Women Knowledge Center
\end{enumerate}

\textbf{Q.144 (24 JUN 2019) | Direct MCQ}
NITI Aayog, GOI, has launched a scheme to develop ecosystem for women entrepreneurship in India. Identify the scheme out of the following :
\begin{enumerate}[label=(\arabic*)]
    \item Women Empowerment Agency
    \item Women Entrepreneurship Platform
    \item Naari Shakti Aayog
    \item Women Knowledge Center
\end{enumerate}

\textbf{Q.109 (DEC 2013) | Direct MCQ}
UNIDO preparatory meeting on the “Role of Women in Industrialisation in Developing Countries” held in February 1978. Identified constraints, which hinder women from participating in industrial activities. The meeting was held at
\begin{enumerate}[label=(\arabic*)]
    \item Vienna
    \item New Delhi
    \item Geneva
    \item Mexico City
\end{enumerate}

\textbf{P-III Q.62 (JUNE 2014) | Direct MCQ}
Indicate the person who is considered as pioneer in the herbal beauty care in India and has achieved unprecedented international acclaim for her application of Ayurveda (the ancient Indian ‘Science of Life’ based on herbal remedies) in cosmetology and trichology :
\begin{enumerate}[label=(\arabic*)]
    \item Sumati Morarji
    \item Shahnaz Hussain
    \item Yamutai Kirloskar
    \item Neena Malhotra
\end{enumerate}

\section{Innovations in Business – Types of Innovations, Creating and Identifying Opportunities, Screening of Business Ideas}

\textbf{Q.5 (JUN 2023) | Multi-Statement}
Which of the statements are true of innovations?
\begin{enumerate}[label=\alph*.]
    \item Innovations are a must to survive in business
    \item Discontinuous innovations lead to failures
    \item Continuous innovations do not disrupt established usage and behaviour patterns
    \item Line extensions are discontinuous innovations
    \item In discontinuous innovations, there is a change not only in technology but also in behavioural patterns of usage and consumption
\end{enumerate}
\begin{enumerate}[label=(\arabic*)]
    \item a, b, d and e only
    \item c, d and e only
    \item b, c and d only
    \item a, c and e only
\end{enumerate}

\textbf{Q.69 (JUN 2023) | Chronological Sequence}
Arrange the steps of novel product design thinking process in a sequential order:
\begin{enumerate}[label=\alph*.]
    \item Define
    \item Test
    \item Prototype
    \item Empathise
    \item Ideate
\end{enumerate}
\begin{enumerate}[label=(\arabic*)]
    \item e, d, a, b, c
    \item a, c, d, e, b
    \item d, a, e, c, b
    \item b, c, d, e, a
\end{enumerate}

\textbf{Q.40 (MAR 2023) | Direct MCQ}
String of activities that moves a product from the raw material stage through manufacturing and distribution, and ultimately to the end user is called:
\begin{enumerate}[label=(\arabic*)]
    \item Business model
    \item Business verticals
    \item Value chain
    \item Business concept
\end{enumerate}

\textbf{Q.41 (MAR 2023) | Direct MCQ}
A metaphor describing the time period in which a firm can realistically enter a new market is called
\begin{enumerate}[label=(\arabic*)]
    \item Opportunity fall
    \item Opportunity gap
    \item Window of opportunity
    \item Opportunity recognition
\end{enumerate}

\textbf{Q.88 (5 SEP 2024) | Direct MCQ}
The process of generating several ideas about a specific topic to generate new business ideas is known as:
\begin{enumerate}[label=(\arabic*)]
    \item Creativity
    \item Discussion
    \item Debate
    \item Brainstorming
\end{enumerate}

\textbf{Q.105 (JUN 2025) | Chronological Sequence}
Arrange the steps of creative process for an idea.
\begin{enumerate}[label=\Alph*.]
    \item Preparation
    \item Incubation
    \item Germination
    \item Verification
    \item Illumination
\end{enumerate}
\begin{enumerate}[label=(\arabic*)]
    \item A, B, C, D, E
    \item C, A, B, E, D
    \item E, D, C, B, A
    \item C, B, A, D, E
\end{enumerate}

\textbf{Q.84 (DEC 2025) | Multi-Statement}
Identify the steps involved in the creative process :
\begin{enumerate}[label=\Alph*.]
    \item Preparation
    \item Opportunity Recognition
    \item Incubation
    \item Insight
    \item Feasibility analysis
\end{enumerate}
\begin{enumerate}[label=(\arabic*)]
    \item A, B and E only
    \item A, C and D only
    \item B, D and E only
    \item A, E and C only
\end{enumerate}

\textbf{Q.120 (DEC 2025) | Match the Following}
Match List - I with List - II.
\begin{tabularx}{\linewidth}{@{} p{0.4\linewidth} X @{}}
    \toprule
    \textbf{List - I Steps to generate creative ideas} & \textbf{List - II Description} \\
    \midrule
    A. Incubation & I. Idea is subject to scrutiny \\
    B. Insight & II. Idea transformed to something of value \\
    C. Evaluation & III. Thinks about a problem \\
    D. Elaboration & IV. Idea is born \\
    \bottomrule
\end{tabularx}
\begin{enumerate}[label=(\arabic*)]
    \item A-III, B-IV, C-I, D-II
    \item A-III, B-I, C-IV, D-II
    \item A-IV, B-III, C-I, D-II
    \item A-IV, B-III, C-II, D-I
\end{enumerate}

\textbf{Q.58 (DEC 2019) | Multi-Statement}
Evaluate the following statements regarding opportunities:
Statement I: All ideas are opportunities.
Statement II: An opportunity is an idea that has the qualities of being attractive, durable and timely.
Which of the following options is correct?
\begin{enumerate}[label=(\arabic*)]
    \item Both the statements are correct
    \item Both the statements are incorrect
    \item Statement I is correct while statement II is incorrect
    \item Statement I is incorrect while statement II is correct
\end{enumerate}

\textbf{Q.90 (4 DEC 2019) | Direct MCQ}
Which one of the following is NOT a source of opportunity identification?
\begin{enumerate}[label=(\arabic*)]
    \item Employees and ex-employees
    \item Dealers and Distributors
    \item Policy Documents
    \item Friends and relatives
\end{enumerate}

\textbf{Q.145 (24 JUN 2019) | Multi-Statement}
Statement I : All ideas are opportunities.
Statement II : An opportunity is an idea that has the qualities of being attractive, durable and timely.
Which of the following options is correct?
\begin{enumerate}[label=(\arabic*)]
    \item Both the statements are correct
    \item Both the statements are incorrect
    \item Statement I is correct while statement II is incorrect
    \item Statement I is incorrect while statement II is correct
\end{enumerate}

\textbf{Q.88 (SEP 2020) | Direct MCQ}
Disruptive innovations in business are those that:
\begin{enumerate}[label=(\arabic*)]
    \item Create a new market by offering value products.
    \item Involve often, new technology or a new business model to offer value to the market.
    \item Displace established market-leading firms, products and alliances.
    \item All of the above.
\end{enumerate}

\textbf{Q.64 (NOV 2021) | Multi-Statement MCQ}
Out of the following, the examples of radical innovations are:
\begin{enumerate}[label=\Alph*.]
    \item Personal Computers
    \item Frozen Yogurt
    \item Disposable diapers
    \item Overnight mail delivery
    \item Microwave popcorn
\end{enumerate}
\begin{enumerate}[label=(\arabic*)]
    \item B, D and E only
    \item A, B and C only
    \item A, C and E only
    \item A, C and D only
\end{enumerate}

\textbf{Q.37 (2022) | Direct MCQ}
Diffusion of innovation theory is associated with:
\begin{enumerate}[label=(\arabic*)]
    \item Peter Drucker
    \item Schumpeter
    \item Everett Rogers
    \item Fredrick Taylor
\end{enumerate}

\textbf{Q.65 (2022) | Multi-Statement}
Innovative small firms are more likely to be in existence in which of the following sectors?
\begin{enumerate}[label=\Alph*.]
    \item Knowledge - based services
    \item Biotechnology
    \item Information technology
    \item Automobiles
    \item Cement
\end{enumerate}
Choose the most appropriate answer:
\begin{enumerate}[label=(\arabic*)]
    \item (D), (E) only
    \item (C), (D), (E) only
    \item (A), (C) only
    \item (B), (E) only
\end{enumerate}

\textbf{Q.75 (2022) | Match the Following}
Match List I with List II:
\begin{tabularx}{\linewidth}{@{} p{0.4\linewidth} X @{}}
    \toprule
    \textbf{List I:} & \textbf{List II:} \\
    \midrule
    (A) Entrepreneurship and innovation & (I) Atal Tinkering lab \\
    (B) Gaps in value chains & (II) Pradhan Mantri Matsya Sampada Yojana \\
    (C) Think Tank & (III) Atal Innovation Mission \\
    (D) Machine learning & (IV) Niti Aayog \\
    \bottomrule
\end{tabularx}
\begin{enumerate}[label=(\arabic*)]
    \item (A)-(II), (B)-(IV), (C)-(I), (D)-(III)
    \item (A)-(I), (B)-(IV), (C)-(III), (D)-(II)
    \item (A)-(IV), (B)-(III), (C)-(I), (D)-(II)
    \item (A)-(III), (B)-(II), (C)-(IV), (D)-(I)
\end{enumerate}

\textbf{Q.119 (DEC 2015) | Multi-Statement}
A business idea becomes a good business opportunity when it covers which four of the following ?
\begin{enumerate}[label=(\alph*)]
    \item Attractiveness
    \item Foreign Collaborations
    \item Political Support
    \item Timeliness
    \item Durability
    \item Value addition to end user
\end{enumerate}
Code :
\begin{enumerate}[label=(\arabic*)]
    \item (a), (b), (d) and (e)
    \item (a), (b), (e) and (f)
    \item (a), (d), (e) and (f)
    \item (b), (c), (e) and (f)
\end{enumerate}

\textbf{Q.41 (AUG 2016 P-III) | Assertion-Reasoning}
Assertion (A) : Business opportunity is an attractive project which motivates the entrepreneur to accept the project for making investment decision.
Reason (R) : A business possibility may take the shape of business opportunity only when it improves social endowments.
Codes:
\begin{enumerate}[label=(\arabic*)]
    \item (A) is correct, but (R) is incorrect.
    \item (A) is incorrect, but (R) is correct.
    \item Both (A) and (R) are correct, but (R) is not the right explanation of (A).
    \item Both (A) and (R) are correct and (R) is the right explanation of (A).
\end{enumerate}

\textbf{Q.43 (NOV 2017 P-II) | Direct MCQ}
Entrepreneurial Judo strategy aims at which one of the following?
\begin{enumerate}[label=(\arabic*)]
    \item Filling holes left by others in the industry.
    \item Developing effectiveness in the industry.
    \item Mobilising finance and effective use to reduce NPAs.
    \item Creating competitive environment in the environment of SMMEs.
\end{enumerate}

\textbf{Q.24 (DEC 2012 P-II) | Direct MCQ}
Which is the basic form of Innovation ?
\begin{enumerate}[label=(\arabic*)]
    \item Introduction of novel production process
    \item Improvement and development of existing process
    \item Improvement and development of existing firm
    \item All of the above
\end{enumerate}

\textbf{Q.57 (JUNE 2013) | Direct MCQ}
Four core themes of innovation are
\begin{enumerate}[label=(\arabic*)]
    \item recognising the opportunity, finding the resources, developing the venture and creating value.
    \item finding the resources, developing the venture, recognizing the opportunity and creating value.
    \item developing the venture, creating value, recognising the opportunity and finding the resources.
    \item creating the value, recognising the opportunity, developing the venture and finding the resources.
\end{enumerate}

\textbf{Q.115 (JUNE 2013) | Direct MCQ}
A sound business opportunity is the result of appropriate interaction between :
\begin{enumerate}[label=(\arabic*)]
    \item needs of the society
    \item capabilities of an entrepreneur
    \item resources available in the environment
    \item all of the above
\end{enumerate}

\textbf{Q.P3-Q60 (SEP 2013) | Match the Following}
Match the following Innovation domains (Paradigm, Position, Product, Process) with their definitions.
\emph{Note: List I, List II and options for this Match the Following question are missing in the provided Markdown.}

\textbf{Q.45 (DEC 2013) | Direct MCQ}
TePP stands for
\begin{enumerate}[label=(\arabic*)]
    \item Technology for Promotion and Production
    \item Trade for Profitability and Productivity
    \item Technopreneur Promotion Programme
    \item None of the above
\end{enumerate}

\section{Business Plan and Feasibility Analysis – Concept and Process of Technical, Market and Financial Analysis}

\textbf{Q.90 (DEC 2023) | Direct MCQ}
Fixing the zero date is important step in establishing a project -
\begin{enumerate}[label=(\arabic*)]
    \item Preliminary steps
    \item Secondary steps
    \item Definition steps
    \item Last steps
\end{enumerate}

\textbf{Q.114 (DEC 2023) | Multi-Statement}
Which of the following are methods of assessment of economic viability of the project?
\begin{enumerate}[label=\Alph*.]
    \item Pay Back Period (PBP)
    \item Return on Investment (ROI)
    \item Net Present Value (NPV)
    \item Benefit Cost Ratio (BCR)
    \item Critical Path Method (CPM)
\end{enumerate}
\begin{enumerate}[label=(\arabic*)]
    \item (A), (B) and (C) Only
    \item (A), (B), (C) and (D) Only
    \item (B), (C), (D) and (E) Only
    \item (A), (C), (D) and (E) Only
\end{enumerate}

\textbf{Q.133 (DEC 2023) | Chronological Sequence}
Arrange the Project Life Cycle Phases in proper order:
\begin{enumerate}[label=\Alph*.]
    \item Concept Phase
    \item Planning and Organising Phase
    \item Definition Phase
    \item Implementation Phase
    \item Project Clean-up Phase
\end{enumerate}
\begin{enumerate}[label=(\arabic*)]
    \item (A), (C), (B), (D), (E) Only
    \item (A), (B), (C), (D), (E) Only
    \item (C), (A), (B), (D), (E) Only
    \item (A), (D), (C), (B), (E) Only
\end{enumerate}

\textbf{Q.79 (MAR 2023) | Chronological Sequence}
Sequence the following steps in the preparation of a versatile and seasoned business plan.
\begin{enumerate}[label=\Alph*.]
    \item Industry and market analysis.
    \item Operations, marketing plan and financing.
    \item Company and management team profiling.
    \item Risk and contingency analysis.
    \item Executive summary.
\end{enumerate}
\begin{enumerate}[label=(\arabic*)]
    \item A, B, D, C, E
    \item C, B, A, D, E
    \item E, C, A, B, D
    \item D, A, C, B, E
\end{enumerate}

\textbf{Q.90 (2 SEP 2024) | Direct MCQ}
What is the primary objective of a feasibility study?
\begin{enumerate}[label=(\arabic*)]
    \item To persuade investors to fund a project
    \item To assess the historical background of a business
    \item To uncover strength and weakness of a business idea
    \item To outline the technical development process
\end{enumerate}

\textbf{Q.87 (5 SEP 2024) | Direct MCQ}
A preliminary evaluation of a business idea for the purpose of determining whether the idea is worth pursuing is known as:
\begin{enumerate}[label=(\arabic*)]
    \item Business plan
    \item Feasibility analysis
    \item Management audit
    \item Business budgeting
\end{enumerate}

\textbf{Q.105 (5 SEP 2024) | Chronological Sequence}
Arrange the following items of a business plan in its proper sequence:
\begin{enumerate}[label=\Alph*.]
    \item Company description
    \item Industry analysis
    \item Market analysis
    \item The Economics of the business
    \item Executive summary
\end{enumerate}
\begin{enumerate}[label=(\arabic*)]
    \item (B), (A), (C), (D), (E)
    \item (E), (B), (A), (C), (D)
    \item (A), (C), (D), (B), (E)
    \item (E), (B), (A), (D), (C)
\end{enumerate}

\textbf{Q.134 (DEC 2025) | Direct MCQ}
The rate at which a company is spending its capital until it reaches profitability is called :
\begin{enumerate}[label=(\arabic*)]
    \item Burn rate
    \item Indifference rate
    \item Break even rate
    \item Spin off rate
\end{enumerate}

\textbf{Q.25 (DEC 2019) | Direct MCQ}
Feasibility analysis is the process of determining if a business idea is viable. It consists of four areas of feasibility analysis. Identify the correct combination.
\begin{enumerate}[label=(\arabic*)]
    \item Product/service feasibility; Industry/target market feasibility; Organizational feasibility; Human resource feasibility
    \item Industry/target market feasibility; Organisational feasibility; Human resource feasibility; Leadership feasibility
    \item Financial feasibility; Product/service feasibility; Industry/target market feasibility; Organizational feasibility
    \item Industry/target market feasibility; Organizational feasibility; Financial feasibility; Leadership feasibility
\end{enumerate}

\textbf{Q.134 (4 DEC 2019) | Multi-Statement}
Marketing feasibility of any idea includes the following aspects :
\begin{enumerate}[label=(\alph*)]
    \item Current and future demand estimates
    \item Market segmentation and identification of target markets
    \item Competition analysis
    \item Market testing
\end{enumerate}
Which of the following option are correct?
\begin{enumerate}[label=(\arabic*)]
    \item (a), (b) and (c) only
    \item (b), (c) and (d) only
    \item (a) and (b) only
    \item (a), (b), (c) and (d)
\end{enumerate}

\textbf{Q.135 (4 DEC 2019) | Multi-Statement}
For marketing the products of MSME, marketing strategy has to be formulated which consists of 4Ps
\begin{enumerate}[label=(\alph*)]
    \item Product
    \item Promotion
    \item Place
    \item Price
\end{enumerate}
Which one of the following sequence is correct?
\begin{enumerate}[label=(\arabic*)]
    \item (a), (b), (c) and (d)
    \item (a), (b), (c) and (d)
    \item (a), (d), (c) and (b)
    \item (a), (c), (d) and (b)
\end{enumerate}

\textbf{Q.149 (24 JUN 2019) | Direct MCQ}
Feasibility analysis is the process of determining if a business idea is viable. It consists of four areas of feasibility analysis. Identify the correct combination.
\begin{enumerate}[label=(\arabic*)]
    \item Product/service feasibility; Industry/target market feasibility; Organizational feasibility; Human resource feasibility
    \item Industry/target market feasibility; Organisational feasibility; Human resource feasibility; Leadership feasibility
    \item Financial feasibility; Product/service feasibility; Industry/target market feasibility; Organizational feasibility
    \item Industry/target market feasibility; Organizational feasibility; Financial feasibility; Leadership feasibility
\end{enumerate}

\textbf{Q.132 (SEP 2020) | Chronological Sequence}
The correct sequence in a business plan can be:
\begin{enumerate}[label=\Alph*.]
    \item Executive summary
    \item Details about the Entrepreneur
    \item Marketing and Economic Details
    \item Financial Details
    \item Human Resources Details
\end{enumerate}
Choose the correct answer from the options given below:
\begin{enumerate}[label=(\arabic*)]
    \item A -> B -> C -> D -> E
    \item B -> C -> E -> D -> A
    \item B -> C -> A -> E -> D
    \item A -> B -> E -> D -> C
\end{enumerate}

\textbf{Q.65 (NOV 2021) | Multi-Statement MCQ}
Technical feasibility analysis will include which of the following?
\begin{enumerate}[label=\Alph*.]
    \item Crucial technical specifications with respect to the design and product safety
    \item Engineering requirements
    \item Product development
    \item Product testing
    \item Plant location
\end{enumerate}
\begin{enumerate}[label=(\arabic*)]
    \item A, B, C and D only
    \item B, C, D and E only
    \item A, B, C, D and E
    \item A, C, D and E only
\end{enumerate}

\textbf{Q.39 (2022) | Direct MCQ}
The first and foremost important feature for a project to be successful is:
\begin{enumerate}[label=(\arabic*)]
    \item It should have a legal support
    \item It should have adequate financing
    \item It should have a market
    \item It should benefit society
\end{enumerate}

\textbf{Q.85 (2022) | Chronological Sequence}
The steps involved in development of a project are given below. Arrange them in proper sequence.
\begin{enumerate}[label=\Alph*.]
    \item Selection of business idea for a detailed analysis from the competing ideas
    \item Project installation and initiation
    \item Feasibility analysis
    \item Identification of investment opportunity
    \item Arrangements for financing
\end{enumerate}
\begin{enumerate}[label=(\arabic*)]
    \item (D), (A), (C), (E), (B)
    \item (A), (C), (E), (D), (B)
    \item (B), (C), (A), (E), (D)
    \item (C), (D), (E), (B), (A)
\end{enumerate}

\textbf{Q.44 (NOV 2017 P-II) | Direct MCQ}
Preparation of a business plan as a pre-requisite to a promotion of a business enterprise, relates to which one of the following?
\begin{enumerate}[label=(\arabic*)]
    \item There is no flaws in the business idea, the fashion, conviction and tenacity of the entrepreneur.
    \item Make arrangement for finance and prevent occurrence of NPAs by business entity.
    \item Ensure that business entity works at break even level.
    \item Make plan to raise finance from market and alternative investment opportunities.
\end{enumerate}

\textbf{Q.25 (DEC 2012 P-III) | Direct MCQ}
A formal document of what the entrepreneur intends to do to sell enough of the firm’s product or service to make a satisfactory profit is called
\begin{enumerate}[label=(\arabic*)]
    \item Long range plan
    \item Strategic plan
    \item Business plan
    \item None of the above
\end{enumerate}

\textbf{Q.46 (JUNE 2012) | Direct MCQ}
Business propositions are to be selected by making an analysis that, how much the local resources will be depleted in the commencement and the course of Business. This is known as
\begin{enumerate}[label=(\arabic*)]
    \item Sensitivity Analysis
    \item Common Property Resource use
    \item Environmental Impact Assessment
    \item None of the above
\end{enumerate}

\textbf{Q.44 (DEC 2013) | Match the Following}
Match the items giving List-I (Market Research, Financial Plan, Ownership, Marketing Plan) into practical functions List-II.
\emph{Note: List I, List II and options for this Match the Following question are missing in the provided Markdown.}

\textbf{Q.P2-Q45 (SEP 2013) | Match the Following}
Match the following Business Planning events (Action planning, Accountability, Working the plan, Reviewing) with descriptions.
\emph{Note: List I, List II and options for this Match the Following question are missing in the provided Markdown.}

\textbf{P-III Q.63 (JUNE 2014) | Direct MCQ}
Which of the following items is not included in project cost estimation of a new business ?
\begin{enumerate}[label=(\arabic*)]
    \item Preference Dividend
    \item Margin money for working capital
    \item Anticipated initial losses
    \item Pre-operative expenses
\end{enumerate}

\section{Micro and Small Scale Industries in India; Role of Government in Promoting SSI}

\textbf{Q.82 (JUN 2023) | Direct MCQ}
Which of the following enterprises is NOT eligible for Udyog Aadhar registration with Ministry of Micro, Small and Medium Enterprises, Government of India?
\begin{enumerate}[label=(\arabic*)]
    \item Micro enterprise in manufacturing sector with an investment in plant and machinery of Rs 20 Lacs
    \item Small enterprise in manufacturing sector with an investment in plant and machinery of Rs 3.5 crores
    \item Small enterprise in services sector with an investment in equipment of Rs 1.50 crores
    \item Medium enterprise in manufacturing sector with an investment in plant and machinery of Rs 15 crores
\end{enumerate}

\textbf{Q.86 (JUN 2023) | Direct MCQ}
“Udyam Assist Platform” developed by the SIDBI caters to which of the following enterprises?
\begin{enumerate}[label=(\arabic*)]
    \item Small Enterprises
    \item Medium Enterprises
    \item Export Oriented Enterprises
    \item Informal Micro Enterprises
\end{enumerate}

\textbf{Q.90 (JUN 2023) | Chronological Sequence}
Arrange the following milk product brands developed by the state co-operatives in ascending order of revenue generator:
\begin{enumerate}[label=\alph*.]
    \item Vita (Haryana)
    \item Amul (Gujarat)
    \item Nandini (Karnataka)
    \item Verka (Punjab)
    \item Milma (Kerala)
\end{enumerate}
\begin{enumerate}[label=(\arabic*)]
    \item e, d, a, b, c
    \item a, e, d, c, b
    \item c, e, a, d, b
    \item d, c, a, e, b
\end{enumerate}

\textbf{Q.53 (MAR 2023) | Multi-Statement}
Which of the following are construed for government intervention to support entrepreneurship and small firms?
\begin{enumerate}[label=\Alph*.]
    \item Stimulate demand and new venture creation in the emerging sectors
    \item Decentralised asset ownership and people participation
    \item Help firms to adopt new technology to avoid market failure
    \item Supporting women participation in business and high-tech firms
    \item Export promotion and currency stabilisation
\end{enumerate}
\begin{enumerate}[label=(\arabic*)]
    \item A, B and C only
    \item C, D and E only
    \item A, C and D only
    \item B, D and E only
\end{enumerate}

\textbf{Q.90 (5 SEP 2024) | Direct MCQ}
Which of the following is not one of the major objectives of the MSME?
\begin{enumerate}[label=(\arabic*)]
    \item Tapping of savings
    \item Export promotion
    \item Utilisation of domestic technology
    \item Rural development
\end{enumerate}

\textbf{Q.124 (5 SEP 2024) | Multi-Statement}
Main components of the National Manufacturing Competitiveness Programme (NMCP) include:
\begin{enumerate}[label=\Alph*.]
    \item Application of lean manufacturing.
    \item Infuse equity funding in MSME.
    \item Management development of MSME.
    \item National campaign for investment in intellectual property.
    \item Provide loans at low interest rates.
\end{enumerate}
\begin{enumerate}[label=(\arabic*)]
    \item (A), (B), (E) Only
    \item (B), (C), (E) Only
    \item (B), (C), (D) Only
    \item (A), (C), (D) Only
\end{enumerate}

\textbf{Q.140 (5 SEP 2024) | Match the Following}
Match List - I with List - II.
\begin{tabularx}{\linewidth}{@{} p{0.4\linewidth} X @{}}
    \toprule
    \textbf{List - I (Name of Scheme)} & \textbf{List - II (Main Objective of Scheme)} \\
    \midrule
    (A) SFURTI & (I) To improve the skills and knowledge of entrepreneurs \\
    (B) ESDP & (II) To create a database of technologies available with various Govt./Pvt. agencies \\
    (C) ASPIRE & (III) To organise traditional industries and artisans into clusters to make them competitive \\
    (D) PMEGP & (IV) To generate employment opportunities in rural as well as urban areas \\
    \bottomrule
\end{tabularx}
\begin{enumerate}[label=(\arabic*)]
    \item (A)-(II), (B)-(I), (C)-(III), (D)-(IV)
    \item (A)-(III), (B)-(I), (C)-(II), (D)-(IV)
    \item (A)-(IV), (B)-(II), (C)-(III), (D)-(I)
    \item (A)-(I), (B)-(III), (C)-(II), (D)-(IV)
\end{enumerate}

\textbf{Q.89 (JUN 2025) | Direct MCQ}
Which of the following institute is not promoting entrepreneurship in India ?
\begin{enumerate}[label=(\arabic*)]
    \item National Institute for Entrepreneurship and Small Business Development (NIESBUD)
    \item Entrepreneurship Development Institute (EDI)
    \item Institute for Design of Electrical Measuring Instruments (IDEMI)
    \item National Small Industries Corporation (NSIC)
\end{enumerate}

\textbf{Q.9 (DEC 2019) | Direct MCQ}
The scheme, which has been launched by the MSME Ministry to assess the capabilities and credit worthiness of the micro and small enterprises (MSEs), is called:
\begin{enumerate}[label=(\arabic*)]
    \item Performance and credit rating scheme
    \item Zero defect zero effect certification
    \item Performance and economic rating scheme
    \item Technology upgradation scheme
\end{enumerate}

\textbf{Q.62 (DEC 2019) | Direct MCQ}
Which of the following were two major criteria for the first official definition of small-scale industry coined in the year 1950?
\begin{enumerate}[label=(\arabic*)]
    \item The size of gross investment in fixed assets and the strength of the work force
    \item The size of gross investment in fixed assets and the annual turnover
    \item The strength of the workforce and the annual turnover
    \item The annual turnover and the annual profit
\end{enumerate}

\textbf{Q.142 (24 JUN 2019) | Direct MCQ}
Which of the following were two major criteria for the first official definition of small scale industry coined, in the year 1950?
\begin{enumerate}[label=(\arabic*)]
    \item the size of gross investment in fixed assets and the strength of the work force
    \item the size of gross investment in fixed assets and the annual turnover
    \item the strength of the workforce and the annual turnover
    \item the annual turnover and the annual profit
\end{enumerate}

\textbf{Q.148 (24 JUN 2019) | Direct MCQ}
The scheme, which has been launched by the MSME Ministry to access the capabilities and creditworthiness of the micro and small enterprises (MSEs), is called :
\begin{enumerate}[label=(\arabic*)]
    \item Performance and credit rating scheme
    \item Zero defect zero effect certification
    \item Performance and economic rating scheme
    \item Technology upgradation scheme
\end{enumerate}

\textbf{Q.89 (SEP 2020) | Direct MCQ}
Marketing Support/Assistance to MSMEs for Bar Code allows:
\begin{enumerate}[label=(\arabic*)]
    \item Reimbursement of 80\% of the one-time registration fee and 80\% of the annual recurring fee for first three years paid by MSMEs to GS1 India for bar coding.
    \item Reimbursement of 75\% of the one-time registration fee only paid by MSMEs to GS1 India for barcoding.
    \item Reimbursement of 75\% of the recurring fee for first three years paid to GS1 by MSMEs for barcoding.
    \item Reimbursement of 75\% of the one-time registration fee and 75\% of the annual recurring fee for first three years paid by MSMEs to GS1 India for barcoding.
\end{enumerate}

\textbf{Q.90 (SEP 2020) | Direct MCQ}
The MSMED Act to cover micro, small and medium Industries in India was introduced in:
\begin{enumerate}[label=(\arabic*)]
    \item June, 2006
    \item June, 2005
    \item June, 2007
    \item August, 2008
\end{enumerate}

\textbf{Q.114 (SEP 2020) | Multi-Statement}
International Cooperation Scheme under marketing promotion scheme to MSME entrepreneurs is available for what type of activities?
\begin{enumerate}[label=\Alph*.]
    \item Participation in International exhibitions, trade fairs.
    \item Giving advertisements in international trade journals.
    \item Participation in buyer-seller meets in foreign countries.
    \item For meeting overseas agents and negotiating business.
    \item For availing overseas grants and aids.
\end{enumerate}
Choose the most appropriate answer from the options given below:
\begin{enumerate}[label=(\arabic*)]
    \item A and C only
    \item A, B and C only
    \item A, C and D only
    \item A and D only
\end{enumerate}

\textbf{Q.40 (2022) | Direct MCQ}
Which scheme on performance and credit rating has been launched by Union MSME Ministry to assess the credit worthiness and capabilities of industries in the sector?
\begin{enumerate}[label=(\arabic*)]
    \item Zero Defect Scheme
    \item Certification Performance and Economy Rating Scheme
    \item Performance and Credit Rating Scheme
    \item Industrial Incentive Scheme
\end{enumerate}

\textbf{Q.125 (DEC 2015) | Direct MCQ}
Which among the following is not a scheme by the MSME ?
\begin{enumerate}[label=(\arabic*)]
    \item Assistance to Training Institutions (ATI)
    \item Credit Linked Capital Subsidy (CLCS)
    \item Credit Linked Assets Subsidy (CLAS)
    \item ISO 9000/ISO 14001 Certification Reimbursement
\end{enumerate}

\textbf{Q.43 (JAN 2017 P-II) | Direct MCQ}
Under technology and quality upgradation support to help manufacturing MSMEs buy energy efficient technologies for production, what percentage of the actual expenditure is provided as financial support by the Government of India?
\begin{enumerate}[label=(\arabic*)]
    \item 75\%
    \item 60\%
    \item 50\%
    \item 25\%
\end{enumerate}

\textbf{Q.34 (JAN 2017 P-III) | Direct MCQ}
The SFURTI programme of the government is to boost production in which one of the following industries?
\begin{enumerate}[label=(\arabic*)]
    \item Auto-components industry
    \item Industries using renewable energy
    \item Khadi and village industry
    \item Construction industry
\end{enumerate}

\textbf{Q.42 (NOV 2017 P-II) | Direct MCQ}
Under the new norms of Micro, Small and Medium Enterprises Development (MSMED) Act 2006, the criteria of investment limit in plant and machinery for MSME falls under which one of the following category?
\begin{enumerate}[label=(\arabic*)]
    \item Micro ₹ 25 lakh, small ₹ 5 crore, Medium ₹ 10 crore
    \item Micro ₹ 20 lakh, small ₹ 2.5 crore, Medium ₹ 5 crore
    \item Micro ₹ 20 lakh, small ₹ 6 crore, Medium ₹ 8 crore
    \item Micro ₹ 25 lakh, small ₹ 10 crore, Medium ₹ 15 crore
\end{enumerate}

\textbf{Q.45 (NOV 2017 P-II) | Direct MCQ}
Federation of Associations of Small Industries of India (FASII) promoted in 1959 is which one of the following type?
\begin{enumerate}[label=(\arabic*)]
    \item Government organisation
    \item Liaisoning organisation
    \item Small Scale unit
    \item SSI Board
\end{enumerate}

\textbf{Q.34 (NOV 2017 P-III) | Direct MCQ}
The support for Entrepreneurial and Managerial Development of Small and Medium Enterprises (incubator) scheme of the National Manufacturing Competitiveness Programme (NMCP) aims to which one of the following?
\begin{enumerate}[label=(\arabic*)]
    \item To provide Startup financial support to new ventures.
    \item To provide financial support to existing MSMEs.
    \item To provide supporting and nourishing business based new ideas in knowledge institutions.
    \item To build awareness on intellectual property rights for the MSMEs
\end{enumerate}

\textbf{Q.35 (NOV 2017 P-III) | Direct MCQ}
The relationship between subsidy and incentive can appropriately be brought out by which one of the following propositions?
\begin{enumerate}[label=(\arabic*)]
    \item Incentives are concessions, bounties offered to borrow money to set up a production activity, while subsidy is a non cash concession given to translate idea into productive action
    \item Incentives and subsidies are cash disbursement only to the intrapreneur
    \item Incentives and subsidies are non-cash motivating factor to promote a new enterprise
    \item Incentives are financial as well as non financial stimuli encouraging to undertake productive activities; while subsidy pertains to a single lumpsum payment given by the government to an entrepreneur for compensating him for the excess cost over administered price of a product or service
\end{enumerate}

\textbf{Q.36 (NOV 2017 P-III) | Direct MCQ}
Industrial parks fall under which one of the following categories?
\begin{enumerate}[label=(\arabic*)]
    \item General industrial clusters
    \item Specialised industrial clusters
    \item Conglomeration of heterogenous group of enterprises under one umbrella
    \item Integrated comprehensive business centres
\end{enumerate}

\textbf{Q.62 (JULY 2018) | Direct MCQ}
In which one of the following years, the Government of India did not bring an Industrial Policy Resolution ?
\begin{enumerate}[label=(\arabic*)]
    \item 1948
    \item 1956
    \item 1965
    \item 1977
\end{enumerate}

\textbf{Q.PA-42 (Archived Rules Set) | Direct MCQ}
The definition of Micro, Small and Medium Enterprises (MSMEs) in India is based on
\begin{enumerate}[label=(\arabic*)]
    \item Total Sales of the Unit
    \item Investment in machines and equipments
    \item Market coverage
    \item Export capacity
\end{enumerate}

\textbf{Q.PA-45 (Archived Rules Set) | Direct MCQ}
In MSMED Act, 2006, the investment limit for Micro, Small and Medium Enterprises have been prescribed as
\begin{enumerate}[label=(\arabic*)]
    \item ₹ 10 lakhs, ₹ 5 crores and ₹ 10 crores
    \item ₹ 5 lakhs, ₹ 1 crore and ₹ 10 crores
    \item ₹ 5 lakhs, ₹ 1 crore and ₹ 5 crores
    \item ₹ 10 lakhs, ₹ 10 crores and ₹ 15 crores
\end{enumerate}

\textbf{Q.PB-41 (Archived Rules Set) | Direct MCQ}
The MSMED Act, 2006 classified enterprises broadly into
\begin{enumerate}[label=(\arabic*)]
    \item Two categories: (i) Manufacturing enterprises and (ii) Service enterprises
    \item \emph{Other options differentiate minor categories, not provided in Markdown.}
\end{enumerate}

\textbf{Q.PB-42 (Archived Rules Set) | Direct MCQ}
Which sector is called as Primary Sector in India?
\begin{enumerate}[label=(\arabic*)]
    \item Industrial sector
    \item Service sector
    \item External sector
    \item Agricultural sector
\end{enumerate}

\textbf{Q.PB-43 (Archived Rules Set) | Direct MCQ}
How many times the census of registered Micro, small and medium Enterprises (MSMEs) have been made by Small Industries Development Organisation (SIDO)?
\begin{enumerate}[label=(\arabic*)]
    \item 1
    \item 2
    \item 3
    \item 4
\end{enumerate}

\textbf{Q.PB-45 (Archived Rules Set) | Direct MCQ}
Industrial Policy of which year contained for the very first time a special thrust on measures for promoting and strengthening small tiny and village enterprises
\begin{enumerate}[label=(\arabic*)]
    \item 1956
    \item 1991
    \item 1948
    \item 1977
\end{enumerate}

\textbf{Q.5 (DEC 2012 P-III) | Direct MCQ}
Micro, Small and Medium Enterprises Development Act, 2006 is related to
\begin{enumerate}[label=(\arabic*)]
    \item Industrial Policy
    \item Investment Limit
    \item Business Opportunities
    \item None of these
\end{enumerate}

\textbf{Q.33 (DEC 2012 P-III) | Direct MCQ}
The Government of India established the Small Industries Development Organisation (SIDO) in the year :
\begin{enumerate}[label=(\arabic*)]
    \item 1951
    \item 1972
    \item 1954
    \item 1982
\end{enumerate}

\textbf{Q.42 (JUNE 2012) | Direct MCQ}
Which of the following organisations is meant for promoting small scale industries at district level ?
\begin{enumerate}[label=(\arabic*)]
    \item SIDBI
    \item DIC
    \item SFCs
    \item NABARD
\end{enumerate}

\textbf{Q.44 (JUNE 2012) | Direct MCQ}
The measure taken by Government to promote small scale industry.
\begin{enumerate}[label=(\arabic*)]
    \item Provision of land
    \item Provision of marketing facilities
    \item Arrangement of credit and raw materials
    \item All the above
\end{enumerate}

\textbf{Q.45 (JUNE 2012) | Direct MCQ}
Small and medium enterprises are
\begin{enumerate}[label=(\arabic*)]
    \item Labour intensive
    \item Capital intensive
    \item Market leader
    \item Industry price determiner
\end{enumerate}

\textbf{Q.111 (DEC 2013) | Direct MCQ}
Small Industries Development Organisation (SIDO) has been set up to formulate, coordinate and monitor the policies and programmes for promotion and development of small scale industries.
\begin{enumerate}[label=(\arabic*)] % Placeholder options for a statement-only question
    \item True
    \item False
    \item Can't be determined
    \item Partially true
\end{enumerate}

\textbf{Q.113 (DEC 2013) | Direct MCQ}
First industrial estate in India was established by SSIB in 1955 at
\begin{enumerate}[label=(\arabic*)]
    \item Okhla in Delhi
    \item \emph{Other options outline other cities, not provided in Markdown.}
\end{enumerate}

\textbf{Q.58 (JUNE 2013) | Direct MCQ}
The main activities of SISI (Small Industries Services Institutes) are
\begin{enumerate}[label=(\arabic*)]
    \item Assistance/Consultancy to prospective Entrepreneurs
    \item Assistance/Consultancy to existing units
    \item (1) and (2)
    \item Assistance/Consultancy and incentives to existing units
\end{enumerate}

\textbf{Q.59 (JUNE 2013) | Chronological Sequence}
Five steps of small business development shall occur in the order of...
\emph{Note: The steps and options for this Chronological Sequence question are missing in the provided Markdown.}

\textbf{Q.60 (JUNE 2013) | Direct MCQ}
Which is not a socio-economic rationale for promoting SSI in India ?
\begin{enumerate}[label=(\arabic*)]
    \item Export Promotion
    \item Employment Generation
    \item Offering competition to large scale industry
    \item Labour intensive
\end{enumerate}

\textbf{Q.119 (JUNE 2013) | Match the Following}
Match training technical institutes with their designated SSI focus domains.
\emph{Note: List I, List II and options for this Match the Following question are missing in the provided Markdown.}

\textbf{Q.44 (DEC 2014) | Direct MCQ}
The number of items which can be exclusively manufactured in the small scale sector are...
\emph{Note: Options for this MCQ are missing.}

\textbf{Q.112 (DEC 2014) | Direct MCQ}
The small scale sector is the ________ largest employer of human resources in the country.
\begin{enumerate}[label=(\arabic*)]
    \item Single
    \item Second
    \item Third
    \item Fourth
\end{enumerate}

\textbf{Q.113 (DEC 2014) | Direct MCQ}
Which of the following is not a promotional measure undertaken by the central government to promote small enterprises ?
\emph{Note: Options for this MCQ are missing.}

\textbf{P-II Q.43 (JUNE 2014) | Direct MCQ}
Which is the apex body in India for formulating the policy in respect of entrepreneurship development ?
\begin{enumerate}[label=(\arabic*)]
    \item Reserve Bank of India
    \item Small Industrial Development Bank of India
    \item National Entrepreneurship Board
    \item Industrial Development Bank of India
\end{enumerate}

\section{Sickness in Small Industries – Reasons and Rehabilitation}

\textbf{Q.38 (MAR 2023) | Direct MCQ}
L.C. Gupta's sickness prediction model primarily focused on which of the following?
\begin{enumerate}[label=(\arabic*)]
    \item Earnings Before Interest and Tax
    \item Sales and Working capital levels
    \item EBIT and operating cash flows
    \item Market value of equity and debt
\end{enumerate}

\textbf{Q.93 (4 DEC 2019) | Multi-Statement}
As per RBI guidelines (2008), an MSME is considered to have reached the stage of incipient sickness if the following events are triggered :
\begin{enumerate}[label=(\alph*)]
    \item When there are cost overruns because of delay in commencement of commercial production by more than 6 months
    \item The enterprise incurs losses for 2 years or cash loss for one year
    \item The capacity utilisation is less than 50 percent of the projected level in term of sales
    \item There is an erosion in the net worth due to accumulated cash losses to the extent of 50 percent of net worth during the previous accounting year
\end{enumerate}
Which one of the following is correct?
\begin{enumerate}[label=(\arabic*)]
    \item (a) and (b) only
    \item (a), (b) and (c) only
    \item (b) and (c) only
    \item (a), (b), (c) and (d)
\end{enumerate}

\textbf{Q.87 (SEP 2020) | Direct MCQ}
Under the Sick Industrial Companies (Special Provision) Act, 1985, a sick Industrial company means:
\begin{enumerate}[label=(\arabic*)]
    \item Being a company registered for not less than 10 years which has at the end of any financial year accumulated losses equal to 50\% of its net worth.
    \item Being a company registered for not less than 5 years which has at the end of any financial year accumulated losses equal to or exceeding its entire net worth.
    \item Being a company which has accumulated losses equal to its entire net worth.
    \item Being a company which immediately after start, has accumulated losses which eat up 50\% of its net worth.
\end{enumerate}

\textbf{Q.64 (2022) | Multi-Statement}
Reason(s) for sickness in small industries in India are:
\begin{enumerate}[label=\Alph*.]
    \item Minimal use of artificial intelligence in business
    \item Inability to cope up with large market size
    \item Lack of innovation and technology focus
    \item Direct and indirect competition from foreign players
    \item Absence of regulatory support
\end{enumerate}
Choose the most appropriate answer:
\begin{enumerate}[label=(\arabic*)]
    \item (B), (D), (E) only
    \item (A), (D), (E) only
    \item (A), (B), (C) only
    \item (B), (C), (D) only
\end{enumerate}

\textbf{Q.117 (DEC 2015) | Multi-Statement}
Which of the following are included as external factors affecting industrial sickness ?
\begin{enumerate}[label=(\alph*)]
    \item Changes in fiscal and monetary policies
    \item Price and distribution
    \item Non-availability of inputs
    \item Industrial relations
    \item Capital market fluctuations
\end{enumerate}
Code :
\begin{enumerate}[label=(\arabic*)]
    \item (a), (b), (c) and (d)
    \item (b), (c), (d) and (e)
    \item Only (a), (b) and (c)
    \item (a), (b), (c), (d) and (e)
\end{enumerate}

\textbf{Q.121 (DEC 2015) | Match the Following}
Match the items of List – I with List – II with regards to the causes of sickness in small scale industry.
\begin{tabularx}{\linewidth}{@{} p{0.4\linewidth} X @{}}
    \toprule
    \textbf{List - I} & \textbf{List - II} \\
    \midrule
    (a) Internal causes of sickness & (i) Own financing by banks \\
    (b) External problems & (ii) Red tapism \\
    (c) Sickness because of government policies & (iii) Faulty project selection \\
    \bottomrule
\end{tabularx}
Codes:
\begin{enumerate}[label=(\arabic*)]
    \item (a)-(i), (b)-(ii), (c)-(iii)
    \item (a)-(i), (b)-(iii), (c)-(ii)
    \item (a)-(ii), (b)-(i), (c)-(iii)
    \item (a)-(iii), (b)-(i), (c)-(ii)
\end{enumerate}

\textbf{Q.122 (DEC 2015) | Assertion-Reasoning}
Assertion (I) : A sick industry is one which is not healthy in terms of its financial management.
Assertion (II) : A unit which has incurred cash losses for one year and it is likely to continue to incur cash loss for current year as well as for following year is considered to be a sick unit.
\begin{enumerate}[label=(\arabic*)]
    \item Assertion (I) only is correct.
    \item Assertion (II) only is correct.
    \item Assertion (I) and (II) both are incorrect.
    \item Assertion (I) and (II) both are correct.
\end{enumerate}

\textbf{Q.123 (DEC 2015) | Direct MCQ}
Which of the following is/are eligible to get the benefit under the scheme of “Rehabilitation of Sick Enterprises” by the Government of India ?
\begin{enumerate}[label=(\arabic*)]
    \item Micro sick unit
    \item Small sick unit
    \item Medium sick unit
    \item All of the above
\end{enumerate}

\textbf{Q.43 (JULY 2016 P-II) | Direct MCQ}
Under the Sick Industrial Companies (Special Provisions) Act, 1985, the public sector companies were covered from which year ?
\begin{enumerate}[label=(\arabic*)]
    \item February, 1994
    \item January, 1990
    \item December, 1991
    \item None of the above
\end{enumerate}

\textbf{Q.45 (JAN 2017 P-II) | Direct MCQ}
Which one of the following is the most insignificant reason for sickness in SSI sector?
\begin{enumerate}[label=(\arabic*)]
    \item Power shortage
    \item Lack of demand
    \item Management problems
    \item Equipment problems
\end{enumerate}

\textbf{Q.60 (DEC 2018) | Assertion-Reasoning}
Given below are two statements, one labelled as Assertion and the other labelled as Reason. Read the statements and choose the correct answer using the code given below.
Assertion "A": Identification and detection of sickness at incipient stage is the first and foremost measure to reduce industrial sickness.
Reason "R": Sickness in small scale industries is not a sudden phenomenon but it is a gradual process taking 5 to 7 years eroding the health of a unit beyond cure.
Codes:
\begin{enumerate}[label=(\arabic*)]
    \item Both "A" and "R" are correct, but "R" is not the right explanation of "A".
    \item "A" is correct, but "R" is incorrect.
    \item Both "A" and "R" are correct and "R" is the right explanation of "A".
    \item Both "A" and "R" are incorrect.
\end{enumerate}

\textbf{Q.57 (JULY 2018) | Direct MCQ}
Timely and adequate assistance and rehabilitation efforts to MSEs should begin on a proactive basis when early signs of sickness are detected. This stage is termed as :
\begin{enumerate}[label=(\arabic*)]
    \item Supporting Stage
    \item Handholding Stage
    \item Sustaining Stage
    \item Recovery Stage
\end{enumerate}

\textbf{Q.59 (JULY 2018) | Assertion-Reasoning}
Assertion (A) : When an industrial unit falls sick, those who depend on it have to face an uncertain future.
Reasoning (R) : The sick units continue to operate below the break-even point and are, thus, forced to depend on external sources for funds of their long-term survival.
Code :
\begin{enumerate}[label=(\arabic*)]
    \item (A) and (R) both are correct ; and (R) is the right explanation of (A).
    \item Both (A) and (R) are correct ; but (R) is not the right explanation of (A).
    \item Both (A) and (R) are incorrect.
    \item (A) is correct but (R) is incorrect.
\end{enumerate}

\textbf{Q.97 (Archived Set 1) | Direct MCQ}
What BIFR stands for :
\begin{enumerate}[label=(\arabic*)]
    \item Bureau of Industrial Financial Reconstruction
    \item Board of Industrial Financial Reconstruction
    \item Bureau of Industrial Financial Re-engineering
    \item Board of Industrial Financial Re-engineering
\end{enumerate}

\textbf{Q.98 (Archived Set 1) | Direct MCQ}
Which one of the following is not the reason for sickness of industries ?
\begin{enumerate}[label=(\arabic*)]
    \item Inadequate working capital
    \item Huge Market
    \item Out dated Technology
    \item Changing tastes of consumers
\end{enumerate}

\textbf{Q.93 (Archived Set 2) | Direct MCQ}
Sick enterprise is referred to
\begin{enumerate}[label=(\arabic*)]
    \item CII
    \item RBI
    \item SIDBI
    \item BIFR
\end{enumerate}

\textbf{Q.96 (Archived Set 2) | Direct MCQ}
Which of the following is the cause of sickness of an enterprise ?
\begin{enumerate}[label=(\arabic*)]
    \item Lack of adequate capital
    \item Lack of demand for the products
    \item Lack of raw material
    \item All of the above
\end{enumerate}

\textbf{Q.PA-43 (Archived Rules Set) | Direct MCQ}
Which one of the following agencies has the power to declare any industrial unit as a potentially sick unit ?
\begin{enumerate}[label=(\arabic*)]
    \item BIFR
    \item SIDBI
    \item IRCI
    \item FICCI
\end{enumerate}

\textbf{Q.42 (DEC 2012 P-II) | Direct MCQ}
Policies related to Revival of Sick Units are framed by
\begin{enumerate}[label=(\arabic*)]
    \item IIFT
    \item CSIR
    \item SEBI
    \item MSME
\end{enumerate}

\textbf{Q.43 (JUNE 2012) | Direct MCQ}
The reason for sickness of small scale industry is
\begin{enumerate}[label=(\arabic*)]
    \item Lack of capital
    \item Lack of market
    \item Severe competition
    \item All the above
\end{enumerate}

\textbf{Q.98 (JUNE 2012) | Direct MCQ}
Sick enterprise is referred to which of the following bodies for rehabilitation ?
\begin{enumerate}[label=(\arabic*)]
    \item Small Industries Development Bank of India (SIDBI)
    \item Small Industries Development Organisation (SIDO)
    \item Board for Industrial and Financial Restructure (BIFR)
    \item National Small Industries Corporation (NSIC)
\end{enumerate}

\textbf{Q.107 (JUNE 2012) | Direct MCQ}
Entrepreneurial failures can be attributed to :
\begin{enumerate}[label=(\arabic*)]
    \item Low quality raw materials
    \item Labour problems
    \item High overhead costs
    \item All of the above
\end{enumerate}

\textbf{Q.117 (JUNE 2013) \& Q.P3-Q58 (SEP 2013) | Direct MCQ}
An industrial unit, according to RBI, is sick if
\begin{enumerate}[label=(\arabic*)]
    \item incurred cash loss in the previous accounting year.
    \item likely to incur loss in the following year.
    \item current ratio is less than 1 : 1 with weak debt: equity ratio.
    \item All of the above.
\end{enumerate}

\textbf{Q.45 (DEC 2014) | Direct MCQ}
Which of the following is an external reason for the sickness of a small business ?
\begin{enumerate}[label=(\arabic*)]
    \item Choice of an idea
    \item Inadequate finance
    \item Volatile business environment
    \item Lack of vertical and horizontal integration
\Hfill
\end{enumerate}

\textbf{P-II Q.45 (JUNE 2014) | Direct MCQ}
Which of the following factors can be regarded as a symptom of industrial sickness ?
\begin{enumerate}[label=(\arabic*)]
    \item Lowering the employee’s morale
    \item Delay and default in the payment of dues
    \item Continuous decrease in the price of its shares
    \item All of the above
\end{enumerate}

\textbf{P-III Q.61 (JUNE 2014) | Direct MCQ}
Which of the following Bill was introduced in the Lok Sabha in 1985, with the objective to make a special provision to secure timely detection of sick and potentially sick companies and for setting up a board of experts for speedy determination of the preventive ameliorative and remedial measures ?
\begin{enumerate}[label=(\arabic*)]
    \item The Industrial Reconstruction Bank of India Bill
    \item NABARD Bill
    \item The Sick Industrial Companies (special provisions) Bill
    \item None of the above
\end{enumerate}

\section{Institutional Finance to Small Industries – Financial Institutions, Commercial Banks, Cooperative Banks, Micro Finance}

\textbf{Q.112 (DEC 2023) | Multi-Statement}
Which of the following are national financial institutions?
\begin{enumerate}[label=\Alph*.]
    \item Industrial Development Bank of India
    \item Industrial Finance Corporation of India
    \item The State Industrial Development Corporation
    \item United Nations Development Programme
    \item State Financial Corporations.
\end{enumerate}
\begin{enumerate}[label=(\arabic*)]
    \item (A), (B) and (C) Only
    \item (B), (C) and (D) Only
    \item (A), (B), (C) and (E) Only
    \item (A), (C), (D) and (E) Only
\end{enumerate}

\textbf{Q.113 (DEC 2023) | Multi-Statement}
Co-operative credit system in India mainly consists of:
\begin{enumerate}[label=\Alph*.]
    \item the short-term agricultural credit institutions.
    \item the long term agricultural credit institutions.
    \item the non-agricultural credit co-operatives.
    \item the short term industrial credit institutions.
\end{enumerate}
\begin{enumerate}[label=(\arabic*)]
    \item (A), (B) Only
    \item (A), (B) and (C) Only
    \item (B), (C) and (D) Only
    \item (A), (C) and (D) Only
\end{enumerate}

\textbf{Q.42 (JUN 2023) | Match the Following}
Match List I with List II:
\begin{tabularx}{\linewidth}{@{} p{0.4\linewidth} X @{}}
    \toprule
    \textbf{LIST I (Financing Concept)} & \textbf{LIST II (Stages of Business Development)} \\
    \midrule
    a. Angel investing & i. Growth and Development stage \\
    b. Sweat equity & ii. Follow-on funding stage \\
    c. Venture capital & iii. Mature stage \\
    d. Initial Public Offer (IPO) & iv. Start-up stage \\
    \bottomrule
\end{tabularx}
\begin{enumerate}[label=(\arabic*)]
    \item a - iii, b - iv, c - i, d - ii
    \item a - iv, b - iii, c - ii, d - i
    \item a - i, b - ii, c - iii, d - iv
    \item a - ii, b - iii, c - i, d - iv
\end{enumerate}

\textbf{Q.54 (MAR 2023) | Multi-Statement}
Indian national strategy for financial inclusion, 2019-24 aims at:
\begin{enumerate}[label=\Alph*.]
    \item Broadening and deepening of financial inclusion
    \item Digitisation and seamless facilitation of financial transactions
    \item Access to formal financial services in an affordable manner.
    \item Women empowerment and strengthening self help group (SHG) ecosystems.
    \item Promoting financial literacy and consumer protection
\end{enumerate}
\begin{enumerate}[label=(\arabic*)]
    \item A, B and C only
    \item C, D and E only
    \item A, B and D only
    \item A, C and E only
\end{enumerate}

\textbf{Q.105 (2 SEP 2024) | Chronological Sequence}
Arrange the financial institutions according to their formation from latest to earliest.
\begin{enumerate}[label=\Alph*.]
    \item Industrial Finance Corporation of India
    \item Regional Rural Bank
    \item National Bank for Agriculture and Rural Development
    \item National Industrial Development Corporation
    \item The Reserve Bank of India.
\end{enumerate}
\begin{enumerate}[label=(\arabic*)]
    \item B, C, A, D, E
    \item C, B, D, A, E
    \item B, D, E, C, A
    \item A, B, C, D, E
\end{enumerate}

\textbf{Q.124 (2 SEP 2024) | Multi-Statement}
Which of the following are operational methodology of Self Help Group?
\begin{enumerate}[label=\Alph*.]
    \item Village meeting of women in villages by MFI
    \item Group formation
    \item Group Training for 5-7 days on procedures and business development skills
    \item Group meeting once in a week for mobilizing saving, disbursing and repaying loans.
    \item Lending money for long term loans
\end{enumerate}
\begin{enumerate}[label=(\arabic*)]
    \item A, B and C Only
    \item B, C and D Only
    \item A, B, C and D Only
    \item B, C, D and E Only
\end{enumerate}

\textbf{Q.125 (2 SEP 2024) | Multi-Statement}
Which of the following Institutions are providing micro finance services to the poors?
\begin{enumerate}[label=\Alph*.]
    \item Conventional weaker section lending banks.
    \item Microfinance Institution.
    \item Self Help Group-bank linkage programme
    \item EXIM banks
\end{enumerate}
\begin{enumerate}[label=(\arabic*)]
    \item A and B Only
    \item B and C Only
    \item C and D Only
    \item A, B and C Only
\end{enumerate}

\textbf{Q.104 (5 SEP 2024) | Chronological Sequence}
Arrange the following stages of Venture Capital Funding in its proper sequence:
\begin{enumerate}[label=\Alph*.]
    \item First \& Second stage of funding
    \item Mezzanine Financing
    \item Start-up Funding
    \item Buyout Funding
    \item Seed Funding
\end{enumerate}
\begin{enumerate}[label=(\arabic*)]
    \item (C), (A), (B), (D), (E)
    \item (E), (C), (A), (B), (D)
    \item (C), (B), (A), (D), (E)
    \item (E), (A), (B), (C), (D)
\end{enumerate}

\textbf{Q.125 (JUN 2025) | Multi-Statement}
Which of the following are not directly financing small scale industries ?
\begin{enumerate}[label=\Alph*.]
    \item Small Industries Development Bank of India (SIDBI)
    \item World Bank
    \item Industrial Development Bank of India (IDBI)
    \item International Monetary Fund (IMF)
    \item National Small Industries Corporation Ltd. (NSIC)
\end{enumerate}
\begin{enumerate}[label=(\arabic*)]
    \item B, C and D Only
    \item A and B Only
    \item A and C Only
    \item A and E Only
\end{enumerate}

\textbf{Q.22 (DEC 2019) | Direct MCQ}
Which out of the following is the principal financial institution for the promotion, financing and development of industries in the small scale sector and to coordinate the functions of other institutions engaged in similar activities?
\begin{enumerate}[label=(\arabic*)]
    \item Infrastructure Development Finance Company Limited
    \item Industrial Finance Corporation of India Limited
    \item Industrial Development Bank of India
    \item Small Industries Development Bank of India (SIDBI)
\end{enumerate}

\textbf{Q.33 (DEC 2019) | Direct MCQ}
The cooperative bank came into existence with the enactment of the Co-operative Credit Societies Act, 1912. Which among the following does not come under the structure of co-operative bank?
\begin{enumerate}[label=(\arabic*)]
    \item State Co-operative Bank
    \item Land Development Bank
    \item Regional Rural Bank
    \item Primary Agricultural Credit Societies
\end{enumerate}

\textbf{Q.146 (24 JUN 2019) | Direct MCQ}
The cooperative bank came into existence with the enactment of the Co-operative Credit Societies Act, 1912. Which among the following does not come under the structure of co-operative bank?
\begin{enumerate}[label=(\arabic*)]
    \item State Co-operative Bank
    \item Land Development Bank
    \item Regional Rural Bank
    \item Primary Agricultural Credit Societies
\end{enumerate}

\textbf{Q.147 (24 JUN 2019) | Multi-Statement}
The following statements relate to Mudra Yojana :
Statement I : Mudra is an NBFC supporting development scheme of micro enterprise sector in the country. It provides refinance support to Bank/MFIs for lending to micro units having loan requirement up to Rs. 10 lakh.
Statement II : Mudra has created different products/schemes for refinancing loans. These interventions have been named 'Shishu', 'Kishore' and 'Tarun' to signify the state of growth/development and funding needs of the beneficiary micro unit/entrepreneur.
Which of the following options is correct?
\begin{enumerate}[label=(\arabic*)]
    \item Statement I is correct, but II is incorrect
    \item Statement II is correct, but I is incorrect
    \item Both the statements I and II are correct
    \item Both the statements I and II are incorrect
\end{enumerate}

\textbf{Q.150 (24 JUN 2019) | Direct MCQ}
Which out of the following is the principal financial institution for the promotion, financing and development of industries in the small scale sector and to coordinate the functions of other institutions engaged in similar activities?
\begin{enumerate}[label=(\arabic*)]
    \item Infrastructure Development Finance Company Limited
    \item Industrial Finance Corporation of India Limited
    \item Industrial Development Bank of India
    \item Small Industries Development Bank of India
\end{enumerate}

\textbf{Q.43 (AUG 2016 P-III) | Match the Following}
Match List - I with List – II.
\begin{tabularx}{\linewidth}{@{} p{0.4\linewidth} X @{}}
    \toprule
    \textbf{List - I (Schemes for SSI)} & \textbf{List - II (Basic function)} \\
    \midrule
    (a) National Equity Fund & (i) Working capital loans and long term loans for fixed capital at one point. \\
    (b) Small Industries Development Fund & (ii) Updated technology support to export oriented units. \\
    (c) Technology development and modernisation fund & (iii) Equity type support to SSI. \\
    (d) Single window scheme & (iv) Refinance assistance to SSI. \\
    \bottomrule
\end{tabularx}
Codes:
\begin{enumerate}[label=(\arabic*)]
    \item (a)-(iii), (b)-(iv), (c)-(ii), (d)-(i)
    \item (a)-(ii), (b)-(i), (c)-(iv), (d)-(iii)
    \item (a)-(i), (b)-(ii), (c)-(iv), (d)-(iii)
    \item (a)-(iv), (b)-(iii), (c)-(i), (d)-(ii)
\end{enumerate}

\textbf{Q.33 (NOV 2017 P-III) | Direct MCQ}
In order to make special allocation of fiscal resources to MSMEs, which one of the following proposals was introduced in the budget of 2015-16?
\begin{enumerate}[label=(\arabic*)]
    \item Micro Units Development Refinance Agency (MUDRA) Bank.
    \item MSME Bank
    \item Swachchha Bharat Abhiyan Vikas Kosh
    \item Self-Employment and Talent Utilisation (SETU) Bank
\end{enumerate}

\textbf{Q.4 (DEC 2018) | Chronological Sequence}
Indicate the code for proper sequencing for the process of venture capital financing from the following:
\begin{enumerate}[label=\arabic*.]
    \item Deal origination
    \item Due diligence
    \item Screening
    \item Deal structuring
    \item Exit plan
\end{enumerate}
Choose the correct answer from the code given below:
\begin{enumerate}[label=(\arabic*)]
    \item (i), (iii), (ii), (iv), (v)
    \item (i), (ii), (iii), (iv), (v)
    \item (i), (ii), (iv), (iii), (v)
    \item (i), (iv), (ii), (iii), (v)
\end{enumerate}

\textbf{Q.47 (Archived Set 1) | Direct MCQ}
Which of the following has been entrusted with the responsibility of developing and supporting small business?
\begin{enumerate}[label=(\arabic*)]
    \item IDBI
    \item RBI
    \item NABARD
    \item SIDBI
\end{enumerate}

\textbf{Q.46 (Archived Set 2) | Direct MCQ}
What is the acronym for SIDBI ?
\begin{enumerate}[label=(\arabic*)]
    \item Short Investment Development Bank of India.
    \item Stock Investment Development Bank of India.
    \item Small Investment Development Bank of India.
    \item Small Industries Development Bank of India.
\end{enumerate}

\textbf{Q.PA-44 (Archived Rules Set) | Direct MCQ}
The headquarters of the SIDBI is in
\begin{enumerate}[label=(\arabic*)]
    \item Mumbai
    \item Lucknow
    \item New Delhi
    \item Hyderabad
\end{enumerate}

\textbf{Q.44 (DEC 2012 P-II) | Direct MCQ}
State Financial Corporation Act encourages in
\begin{enumerate}[label=(\arabic*)]
    \item Establishing Industrial Estates
    \item Establishing small and medium size industries
    \item Establishing Agricultural Farms
    \item All of the above
\end{enumerate}

\textbf{Q.48 (JUNE 2012) | Direct MCQ}
Funding through Micro finance to Small and Micro Enterprises is not successful due to
\begin{enumerate}[label=(\arabic*)]
    \item High rate of interest
    \item Out sourcing of method of providing Micro Finance
    \item (1) \& (2)
    \item None of the above
\end{enumerate}

\textbf{Q.62 (JUNE 2012) | Match the Following}
Match the following Financial Institutions (IDBI, SIDBI, NSIC, IFCI) with their corresponding Year of Establishment.
\emph{Note: List I, List II and options for this Match the Following question are missing in the provided Markdown.}

\textbf{Q.P3-Q57 (SEP 2013) | Direct MCQ}
The main function of venture capitalist is to
\begin{enumerate}[label=(\arabic*)]
    \item Provide funds to the steps necessary to establish the commercial viability of a new product.
    \item \emph{Other options detail various funding stages like marketing, production, or research alone, not provided in Markdown.}
\end{enumerate}

\textbf{Q.P3-Q59 (SEP 2013) | Direct MCQ}
SIDBI has been entrusted with the responsibility of
\begin{enumerate}[label=(\arabic*)]
    \item Providing financial assistance to all sectors.
    \item Developing medium and large scale industry.
    \item Providing incentives to medium scale industry.
    \item Developing and supporting small business.
\end{enumerate}

\textbf{P-II Q.42 (JUNE 2014) | Direct MCQ}
________ is considered as a creative capital which performs economic functions different from other investment vehicles, which primarily serve as the expansion capital.
\begin{enumerate}[label=(\arabic*)]
    \item Equity Capital
    \item Fixed Capital
    \item Venture Capital
    \item Share Capital
\end{enumerate}

\textbf{P-III Q.64 (JUNE 2014) | Direct MCQ}
Which of the following is the main function of State Financial Corporations ?
\begin{enumerate}[label=(\arabic*)]
    \item Provide short-term finance to large-sized industrial units
    \item Provide long-term finance to small and medium sized industrial units.
    \item Undertake variety of promotional activities, like preparation of feasibility reports, conducting industrial potential surveys, etc.
    \item Develop industrial estates for providing land to a building Entrepreneurs.
\end{enumerate}

\section{Information Technology \& Analytics (Interdisciplinary questions from combined exam models)}

\emph{Note: The following questions were included in the DEC 2025 paper alongside Entrepreneurship questions, representing an interdisciplinary integration with Information Technology and Analytics. Case text for Reading Comprehension/Case-Based MCQs is not provided.}

\textbf{Q.55 (DEC 2025) | Chronological Sequence}
Based on the Data Mining Process used in Business Analytics arrange the following in sequence :
\begin{enumerate}[label=\Alph*.]
    \item Pre-Processing the data selecting attributes of interest and checking the outliers
    \item Assess the degree to which the selected model make company understood variability/similarity in behaviour and generate segments
    \item Understanding "What are the common characteristics of the customers, that the company has lost to its competitors"
    \item Identify the spending behaviour of shoppers understanding the most relevant variables and use of descriptive statistics.
    \item Modeling technique are selected and applied on dataset.
\end{enumerate}
\begin{enumerate}[label=(\arabic*)]
    \item B, E, C, D, A
    \item C, E, B, A, D
    \item D, E, B, A, C
    \item C, D, A, E, B
\end{enumerate}

\textbf{Q.57 (DEC 2025) | Direct MCQ}
Which of the following is not the type of Data Warehouse ?
\begin{enumerate}[label=(\arabic*)]
    \item Data Marts (DMs)
    \item Operational Data Stores (ODS)
    \item Enterprise Data Warehouses (EDW)
    \item Data Integration Houses (DIH)
\end{enumerate}

\textbf{Q.76 (DEC 2025) | Multi-Statement}
The "Vs" that define the Big Data are :
\begin{enumerate}[label=\Alph*.]
    \item Volume
    \item Variety
    \item Velocity
    \item Validity
    \item Veracity
